\documentclass[11pt,a4paper]{article}
\usepackage[utf8]{inputenc}
\usepackage[T1]{fontenc}
\renewcommand{\contentsname}{Indice}
\renewcommand{\partname}{Parte}
\renewcommand{\figurename}{Figura}
\renewcommand{\tablename}{Tabella}
\usepackage[margin=2.1cm]{geometry}
\usepackage{amsmath,amssymb}
\usepackage{enumitem}
\usepackage{xcolor}
\usepackage{longtable}
\usepackage{array}
\usepackage[colorlinks=true,linkcolor=blue!55!black,urlcolor=blue!55!black]{hyperref}

\usepackage{framed}
\setlength{\parindent}{0pt}
\setlength{\parskip}{5pt}
\newcommand{\qid}[1]{\textbf{[#1]}}
\newcommand{\OO}[1]{O\!\left(#1\right)}
\newcommand{\TT}[1]{\Theta\!\left(#1\right)}
\newcommand{\Om}[1]{\Omega\!\left(#1\right)}
% --- riquadro col testo dell'esercizio ---
\definecolor{shadecolor}{RGB}{235,242,250}
\newenvironment{domanda}%
{\begin{shaded}\small\textbf{Testo dell'esercizio.} \itshape}%
{\end{shaded}}
% --- evidenziazione dei punti chiave ---
\newcommand{\kw}[1]{\ifmmode{\color{red!60!black}#1}\else\textcolor{red!60!black}{\textbf{#1}}\fi}
% --- marcatore per gli argomenti usciti in 3+ prove ---
\newcommand{\hot}{\,\textcolor{orange!85!black}{$\bigstar$}}

\title{\textbf{Algorithm Engineering}\\[2mm]
\Large Esercizi di Teoria --- testi e risoluzioni dettagliate\\[1mm]
\normalsize Tutte le richieste di teoria di appelli, midterm, final term e orali 2020--2025}
\author{}
\date{Versione riveduta e completata --- luglio 2026}

\begin{document}
\maketitle
\vspace{-1.2cm}

\section*{Come usare questo documento}

\textbf{Parte I}: i testi delle domande di teoria, organizzati per argomento. Ogni testo ha un
identificatore (es.\ \qid{5.1}) e accanto sono indicate le prove d'esame in cui \`e uscito.
\textbf{Parte II}: le risoluzioni, con dimostrazioni complete e commenti, sviluppate a partire
dalla dispensa del corso e da \emph{Pearls of Algorithm Engineering} (Ferragina, 2023).
La numerazione delle risoluzioni rispecchia quella dei testi.

\textbf{Copertura.} Questa versione \`e stata verificata contro \emph{tutte} le prove presenti
nella cartella \texttt{esercizi/} (35 prove con testo leggibile, dal 10/01/2020 al 15/07/2025,
inclusi i .doc del 2020 e del 2022 e l'orale del 25/07/2022). Rispetto alla versione precedente
sono stati aggiunti gli item mancanti:
\qid{1.5} sorting vs permuting nel modello RAM e I/O (orale 25/07/22);
\qid{3.3} sampling con $n$ noto e prova di correttezza (midterm 06/11/23);
\qid{5.9} prova che l'altezza media del Treap \`e $\OO{\log n}$ (07/02/20, orale 25/07/22);
\qid{7.7} algoritmo lineare di Kasai per l'array LCP (orale 25/07/22, domanda ``da lode'');
\qid{8.3} come ricavare $n$ dagli array $L,H$ di Elias--Fano (orale 25/07/22);
\qid{9.3} bound sull'errore di troncamento di un numero binario (10/01/20);
e la nuova sezione 11 con le tre \emph{domande di progettazione} (07/02/20, 02/02/22, 11/11/24).
Le convenzioni (bit alti di Elias--Fano, tabelle FC/SYMB, Wheeler code) sono state allineate a
quelle effettivamente usate nei compiti del docente.

\textbf{Legenda.} Nella Parte II ogni risoluzione \`e preceduta dal \emph{testo dell'esercizio}
(riquadro azzurro) con le date in cui \`e uscito; i passaggi e i risultati da ricordare sono
\kw{evidenziati in rosso} o incorniciati. Il simbolo \hot{} marca gli argomenti usciti in
\emph{tre o pi\`u} prove: da sapere alla perfezione.

\textbf{Nota sulle soluzioni ufficiali.} Le soluzioni ufficiali esistono solo come scansioni
manoscritte della parte di esercizi pratici; per la teoria non esistono soluzioni ufficiali in
formato testo, quindi le risoluzioni qui riportate sono ricostruite dalle fonti del corso.

\subsection*{Tabella di copertura: prova per prova}

\begin{longtable}{@{}p{2.6cm}p{2.2cm}p{10.2cm}@{}}
\textbf{Data} & \textbf{Tipo} & \textbf{Domande di teoria e riferimenti}\\
\hline
\endhead
10/01/20 & appello & cammini cuckoo $c^{-L}/r$ \qid{5.3}; spazio Elias--Fano \qid{8.1}; errore di troncamento binario \qid{9.3}\\
07/02/20 & appello & design I/O max-satellite \qid{11.1}; bound sorting $D$ dischi con prova e confronto con l'ottimo \qid{1.3}\qid{1.4}; profondit\`a media Treap \qid{5.9}\\
12/06/20 & appello & suffix tree: sottostringa con $\ge L$ occorrenze \qid{7.6}\\
10/07/20 & appello & compacted trie con fan-out hash per ordinare stringhe \qid{6.5}\\
11/11/20 & midterm & Snow Plow con prob.\ heap $3/4$ \qid{2.1}; collisioni attese con hash universale \qid{5.2}; Bloom: $p_0=\tfrac12$ con $k$ ottimo \qid{5.4}\\
07/09/21 & appello & costruzione SA via qsort, e caso LCP $\le c$ \qid{7.3}; perch\'e l'aritmetico parte dall'elemento centrale \qid{9.6}\\
14/01/22 & final term & intersezione di liste Elias--Fano con \texttt{Access}/\texttt{GEQ} \qid{4.4}\\
02/02/22 & appello & design: finestra di $k$ posizioni con $q$ occorrenze via SA \qid{11.2}\\
25/07/22 & orale & sorting/permuting in RAM e I/O \qid{1.5}; lower bound stringhe + Multikey QS \qid{2.2}\qid{2.3}; Treap: definizione e altezza media \qid{5.8}\qid{5.9}; ricavare $n$ da Elias--Fano \qid{8.3}; LCP in tempo lineare (lode) \qid{7.7}\\
16/01/23 & final term & LSD RadixSort \qid{2.4}; hash universale \qid{5.1}; Rank \qid{10.1}; def.\ SA+LCP \qid{7.1}\\
10/02/23 & appello & teorema cuckoo \qid{5.3}; Multikey QS + lower bound \qid{2.2}\qid{2.3}; upper bound I/O sort/permute/striping \qid{1.2}; Treap: definizione e split \qid{5.8}\\
05/06/23 & appello & lower bound + Multikey QS + RadixSort ottimo \qid{2.2}\qid{2.3}\qid{2.4}; Rank + spazio \qid{10.1}; Bloom su 2 server \qid{5.6}\\
24/07/23 & appello & SA via qsort e I/O \qid{7.3}; ricerca lessicografica con Patricia \qid{6.1}; Bloom: definizione + falso positivo \qid{5.5}\\
07/09/23 & appello & Multikey QS \qid{2.3}; hash universale \qid{5.1}; SA/LCP + ricerca + prova \qid{7.2}\\
02/11/23 & appello & Reservoir + correttezza \qid{3.2}; hash universale \qid{5.1}; Rank \qid{10.1}; def.\ SA+LCP \qid{7.1}\\
06/11/23 & midterm & sampling con $n$ noto + correttezza ($m=1$) \qid{3.3}; universale + chaining $\OO{1}$ \qid{5.2}\\
12/12/23 & final term & upper bound Arithmetic \qid{9.2}; spazio della struttura di Rank \qid{10.2}\\
15/01/24 & appello & Snow Plow $2M$ e $p=\tfrac14$ \qid{2.1}; COUNT con SA + retrieval su disco \qid{7.4}; doubling \qid{4.3}\\
09/02/24 & appello & MergeSort multi-way I/O + disk striping \qid{1.1}\qid{1.3}; Canonical Huffman \qid{9.1}; intersezione two-level \qid{4.1}\\
27/05/24 & appello & Bloom \qid{5.4}; spazio Elias--Fano \qid{8.1}; MergeSort I/O \qid{1.1}; storage di un dizionario di stringhe \qid{6.4}\\
21/06/24 & appello & Snow Plow con $p=\tfrac18$ \qid{2.1}; two-level scan + prova \qid{4.1}; Canonical Huffman \qid{9.1}\\
11/07/24 & appello & Reservoir ($m<n/2$) + correttezza \qid{3.2}; LSD su $b$ bit \qid{2.4}; mutual partition \qid{4.2}; Access su Elias--Fano \qid{8.2}\\
11/11/24 & midterm & altezza Skip List \qid{5.7}; errore del Bloom \qid{5.4}; two-level intersection \qid{4.1}; design range di frequenze \qid{11.3}\\
09/12/24 & final term & upper bound Arithmetic \qid{9.2}\\
10/01/25 & appello & sorting con $D$ dischi vs ottimo \qid{1.4}; Bounded QuickSort \qid{2.5}; hash universale \qid{5.1}; LSD + correttezza \qid{2.4}; Rank $\OO{1}$/$o(m)$ \qid{10.2}\\
04/02/25 & appello & MergeSort I/O con 1 e $D$ dischi \qid{1.1}\qid{1.3}; Bounded QuickSort \qid{2.5}; Skip List \qid{5.7}; Rank \qid{10.2}; spazio Elias--Fano \qid{8.1}\\
06/03/25 & appello & pseudocodice Reservoir \qid{3.1}; bit di Huffman vs Arithmetic + quando \qid{9.4}; Compacted vs Patricia \qid{6.2}\\
03/06/25 & appello & pseudocodice Reservoir \qid{3.1}; Huffman vs Arithmetic \qid{9.4}; teorema cuckoo \qid{5.3}; disk striping \qid{1.3}; Compacted vs Patricia nel two-level \qid{6.3}\\
23/06/25 & appello & LSD + correttezza \qid{2.4}; Bounded QuickSort \qid{2.5}; definizione Treap \qid{5.8}; invarianza per permutazione dell'Arithmetic \qid{9.5}\\
15/07/25 & appello & errore del Bloom \qid{5.4}; two-level intersection \qid{4.1}; ricerca di $P$ in SA \qid{7.5}; Canonical Huffman \qid{9.1}\\
\hline
\end{longtable}

Prove \emph{senza} domande di teoria (solo esercizi di simulazione): 03/09/20, 16/12/20,
20/01/21, 10/02/21 (a parte il commento di complessit\`a del COUNT, coperto da \qid{7.4}),
09/04/21, 08/06/21, 28/10/21, 08/11/21, 15/12/21, 11/04/22, 13/06/22, 04/07/22, 25/07/22 (scritto).

\tableofcontents
\newpage

\part{Testi degli esercizi di teoria}
\section{Modelli e complessit\`a I/O}

\begin{itemize}[leftmargin=1.4em]
\item \qid{1.1}\hot Enuncia la complessit\`a I/O del MergeSort multi-way in funzione di $M$, $B$, $N$.
\emph{(09/02/24, 27/05/24, 04/02/25)}
\item \qid{1.2} Enuncia gli upper bound per: ordinare $n$ item atomici in un modello a due livelli con
memoria interna $M$, pagina $B$ e 1 disco; permutarli nello stesso modello con 1 disco; ordinarli con
$D$ dischi usando il disk striping. \emph{(10/02/23)}
\item \qid{1.3}\hot Descrivi la tecnica del disk striping ed enuncia il bound I/O del MergeSort quando lo
si usa per ordinare $n$ item in una memoria a due livelli con $D$ dischi, memoria $M$, pagina $B$.
\emph{(07/02/20 con prova, 09/02/24, 04/02/25, 03/06/25)}
\item \qid{1.4} Enuncia la complessit\`a I/O di ordinare $n$ item con memoria $M$, pagina $B$ e $D$
dischi; confrontala con il bound ottimo e discuti le differenze. \emph{(07/02/20, 10/01/25)}
\item \qid{1.5} Dati $n$ item atomici, scrivi la complessit\`a in tempo dei problemi di sorting e di
permuting nel modello RAM, poi la loro complessit\`a I/O nel modello a due livelli (pagina $B$,
memoria $M$); commenta se i due problemi sono computazionalmente equivalenti in senso asintotico e
per quali valori dei parametri. \emph{(orale 25/07/22)}
\end{itemize}

\section{Ordinamento}

\begin{itemize}[leftmargin=1.4em]
\item \qid{2.1}\hot Dimostra che la lunghezza attesa della sequenza ordinata (run) prodotta da Snow Plow
\`e $2M$. Varianti chieste: quanto vale se la probabilit\`a di finire nell'\emph{unsorted bucket} \`e
$\tfrac14$ invece di $\tfrac12$? \emph{(15/01/24)} Se \`e $\tfrac18$? \emph{(21/06/24)} Se la
probabilit\`a di andare nello \emph{heap} \`e $\tfrac34$? \emph{(11/11/20)}
\item \qid{2.2}\hot Enuncia e dimostra il lower bound (comparison-based) al tempo per ordinare $n$
stringhe di lunghezza totale $N$ su alfabeto di taglia $\sigma$. \emph{(10/02/23, 05/06/23, orale 25/07/22)}
\item \qid{2.3}\hot Scrivi lo pseudocodice del Multikey QuickSort su $n$ stringhe di lunghezza totale $N$;
enuncia e dimostra la sua complessit\`a temporale e mostra che \`e time-optimal rispetto al lower
bound di \qid{2.2}. \emph{(10/02/23, 05/06/23, 07/09/23, orale 25/07/22)}
\item \qid{2.4}\hot Abbozza l'algoritmo LSD RadixSort per $n$ chiavi di $b$ bit ciascuna; enuncia e
dimostra la sua complessit\`a temporale; dimostra la correttezza dell'ordinamento prodotto (caso
stringhe binarie). Variante: complessit\`a del RadixSort ottimo su $n$ stringhe binarie di lunghezza
totale $N$. \emph{(16/01/23, 05/06/23, 11/07/24, 10/01/25, 23/06/25)}
\item \qid{2.5}\hot Scrivi e commenta lo pseudocodice del Bounded QuickSort, cio\`e la variante che
garantisce uno stack di taglia $\OO{\log n}$ (equivalentemente: $\OO{\log n}$ chiamate ricorsive
annidate per qualunque input). \emph{(10/01/25, 04/02/25, 23/06/25)}
\end{itemize}

\section{Random Sampling}

\begin{itemize}[leftmargin=1.4em]
\item \qid{3.1}\hot Scrivi (solo) lo pseudocodice del Reservoir Sampling. \emph{(06/03/25, 03/06/25)}
\item \qid{3.2}\hot Data una sequenza $S$ di $n$ item ($n$ ignoto) e un intero $m$ (eventualmente
$m<n/2$): mostra lo pseudocodice del Reservoir Sampling e dimostrane la correttezza, cio\`e che ogni
item \`e campionato con probabilit\`a $m/n$ (campionamento uniforme). \emph{(02/11/23, 11/07/24)}
\item \qid{3.3} Data una sequenza $S$ di $n$ item con $n$ \emph{noto} e un intero $m$, scrivi lo
pseudocodice dell'algoritmo di sampling che estrae $m$ item uniformemente a caso, e dimostrane la
correttezza (probabilit\`a $m/n$) nel caso $m=1$. \emph{(06/11/23; simulazioni in 13/06/22, 10/02/23, 07/09/23)}
\end{itemize}

\section{Set Intersection}

\begin{itemize}[leftmargin=1.4em]
\item \qid{4.1}\hot Date due liste ordinate $L_1$, $L_2$ di lunghezze $n$ e $m$: descrivi l'algoritmo
\emph{two-level} (detto anche two-level scan / two-level memory / two-level storage) per calcolarne
l'intersezione; enuncia e dimostra la sua complessit\`a. \emph{(09/02/24, 21/06/24, 11/11/24, 15/07/25)}
\item \qid{4.2} Descrivi l'algoritmo \emph{mutual partition} per l'intersezione di $L_1$, $L_2$,
con la sua complessit\`a. \emph{(11/07/24)}
\item \qid{4.3} Descrivi il \emph{doubling algorithm} (binary search con salti esponenziali,
``galloping'') per l'intersezione; enuncia e dimostra la sua complessit\`a. \emph{(15/01/24)}
\item \qid{4.4} Scrivi lo pseudocodice per intersecare due liste ordinate $L_1$, $L_2$ (di $n_1$ e
$n_2$ elementi) codificate con Elias--Fano, usando come primitive $\texttt{Access}(L,i)$ e
$\texttt{GEQ}(L,x)$; valuta e commenta la complessit\`a. \emph{(14/01/22)}
\end{itemize}

\section{Hashing, Bloom Filter, Skip List, Treap}

\begin{itemize}[leftmargin=1.4em]
\item \qid{5.1}\hot Definisci una classe di funzioni hash universali; fornisci un esempio e dimostrane
l'universalit\`a. \emph{(16/01/23, 07/09/23, 02/11/23, 10/01/25)}
\item \qid{5.2} Definisci una classe universale e mostra che, usata nell'hashing con chaining,
ottiene lunghezza attesa $\OO{1}$ per catena quando $m=\OO{n}$; da qui il calcolo del numero atteso
di collisioni. \emph{(06/11/23; esercizi collegati in 10/07/20 e 11/11/20)}
\item \qid{5.3}\hot Enuncia e dimostra il teorema sul cuckoo graph: per ogni coppia di posizioni $i,j$
della tabella cuckoo, la probabilit\`a che il cammino minimo $i\to j$ abbia lunghezza $L$ \`e al
pi\`u $c^{-L}/r$, dove $r$ \`e il numero di posizioni della tabella, $n$ le chiavi memorizzate e
$r\ge 2cn$. Commenta l'impatto sul load factor. \emph{(10/01/20, 10/02/23, 03/06/25)}
\item \qid{5.4}\hot Calcola la probabilit\`a che una posizione del bit-array del Bloom Filter (taglia
$m$, $n$ chiavi inserite) sia 0, e deriva/dimostra la probabilit\`a d'errore complessiva. Variante:
dimostra che tale probabilit\`a \`e $\tfrac12$ quando il numero di funzioni hash \`e quello ottimo.
\emph{(11/11/20, 27/05/24, 11/11/24, 15/07/25)}
\item \qid{5.5}\hot Definisci un Bloom Filter su un dizionario di $n$ chiavi e le operazioni supportate;
enuncia e dimostra il risultato sul falso positivo. \emph{(24/07/23)}
\item \qid{5.6} Dati due insiemi $A$, $B$ memorizzati su due server, mostra come calcolarne
l'intersezione usando un Bloom Filter e un solo round di comunicazione, ammettendo errori.
\emph{(05/06/23)}
\item \qid{5.7} Dimostra che l'altezza (profondit\`a) media di una Skip List con $n$ nodi \`e
$\OO{\log n}$. \emph{(11/11/24, 04/02/25)}
\item \qid{5.8}\hot Definisci la struttura Treap su $n$ coppie $\langle$chiave, priorit\`a$\rangle$ con
priorit\`a minima nella radice; descrivi le operazioni supportate e come si implementa lo split su
una chiave $K$ non presente. \emph{(10/02/23, 23/06/25, orale 25/07/22)}
\item \qid{5.9} Enuncia e dimostra che la profondit\`a media di un Treap con priorit\`a scelte
uniformemente a caso \`e logaritmica nel numero di chiavi. \emph{(07/02/20, orale 25/07/22)}
\end{itemize}

\section{Prefix Search}

\begin{itemize}[leftmargin=1.4em]
\item \qid{6.1} Dato un Patricia Trie su un dizionario $S$ di $n$ stringhe, descrivi come cercare la
posizione lessicografica di $P[1,p]$ in $S$; enuncia la complessit\`a in tempo e in I/O.
\emph{(24/07/23)}
\item \qid{6.2} Descrivi la differenza fra Compacted Trie e Patricia Trie in termini di spazio e
tempo di query, in funzione di $n$, della lunghezza totale $N$ e della lunghezza del pattern $P$.
\emph{(06/03/25)}
\item \qid{6.3} Qual \`e la differenza principale fra Compacted Trie e Patricia Trie, e come \`e
stata sfruttata per costruire un indice \emph{two-level} efficiente per un insieme di stringhe?
\emph{(03/06/25)}
\item \qid{6.4} Dato un dizionario $D$ di $n$ stringhe di lunghezza variabile e lunghezza totale
$N$, discuti almeno 3 soluzioni di memorizzazione, commentando spazio e costo tempo/I/O di
$\texttt{Access}(i)$ (recupero della $i$-esima stringa). \emph{(27/05/24)}
\item \qid{6.5} Descrivi come usare un compacted trie (con fan-out dei nodi in hash table) per
ordinare $n$ stringhe di lunghezza totale $L$; commenta tempo e spazio. \emph{(10/07/20)}
\end{itemize}

\section{Substring Search}

\begin{itemize}[leftmargin=1.4em]
\item \qid{7.1} Definisci formalmente il suffix array $SA$ di un testo $T[1,n]$ e il corrispondente
array LCP. \emph{(16/01/23, 02/11/23)}
\item \qid{7.2}\hot Come \qid{7.1}, e inoltre: descrivi come cercare $P[1,p]$ in $SA$; enuncia e dimostra
la complessit\`a della ricerca. \emph{(07/09/23)}
\item \qid{7.3} Definisci $SA$; scrivi lo pseudocodice per costruirlo via qsort; enuncia e commenta
la complessit\`a worst-case (tempo e I/O) di questa costruzione, e il caso in cui il massimo LCP fra
i suffissi \`e limitato da una costante $c$. \emph{(07/09/21, 24/07/23)}
\item \qid{7.4}\hot Mostra come CONTARE in $T[1,n]$ tutte le occorrenze di $P[1,p]$ con un Suffix Array
in memoria; enuncia e dimostra la complessit\`a; qual \`e il costo I/O del RETRIEVAL delle posizioni
se $SA$ \`e su disco? \emph{(10/02/21, 15/01/24)}
\item \qid{7.5}\hot Descrivi come cercare $P[1,p]$ in un $SA$ costruito su $T[1,n]$ e commenta la
complessit\`a in funzione di $p$ e $n$. \emph{(15/07/25)}
\item \qid{7.6} Dato il suffix tree di $T[1,n]$, come si calcola la sottostringa pi\`u lunga che
occorre almeno $L$ volte? Complessit\`a in funzione di $n$. \emph{(12/06/20)}
\item \qid{7.7} Scrivi lo pseudocodice dell'algoritmo lineare (Kasai et al.) che calcola l'array LCP
dato $T$ e $SA$, e dimostra che \`e $\OO{n}$. \emph{(orale 25/07/22, domanda ``da lode''; simulazioni
in 11/11/24, 06/03/25, 15/07/25)}
\end{itemize}

\section{Integer Coding}

\begin{itemize}[leftmargin=1.4em]
\item \qid{8.1}\hot Data una sequenza monotona di $n$ interi non negativi $<u$, enuncia e dimostra lo
spazio (in bit) della codifica di Elias--Fano. \emph{(10/01/20, 27/05/24, 04/02/25)}
\item \qid{8.2} Scrivi l'algoritmo per eseguire $\texttt{Access}(i)$ su una sequenza codificata con
Elias--Fano; cenno a $\texttt{NextGEQ}(x)$. \emph{(11/07/24; simulazioni in molti appelli)}
\item \qid{8.3} Dati gli array $L$ e $H$ di una codifica Elias--Fano, mostra come ricavare il valore
di $n$ (numero di interi codificati), motivando la risposta. \emph{(orale 25/07/22; \`e il ``hint''
ricorrente degli esercizi: 11/11/20, 12/12/23, 09/12/24, 06/03/25, 23/06/25)}
\item \qid{8.4} (Riferimento per gli esercizi di calcolo) Numero di bit emessi dai codici unario,
$\gamma$, $\delta$, Rice e $(s,c)$-dense su un intero $x$. \emph{(vari appelli)}
\end{itemize}

\section{Statistical Coding}

\begin{itemize}[leftmargin=1.4em]
\item \qid{9.1}\hot Descrivi perch\'e \`e stato introdotto il Canonical Huffman; specifica quali
strutture dati tiene nel preambolo del file compresso; scrivi lo pseudocodice per decomprimere un
simbolo. \emph{(09/02/24, 21/06/24, 15/07/25)}
\item \qid{9.2} Dimostra l'upper bound in bit della codifica aritmetica, in funzione dell'entropia e
della lunghezza del testo in input. \emph{(12/12/23, 09/12/24)}
\item \qid{9.3} Dimostra un upper bound alla variazione indotta dal troncamento a $k$ cifre di un
numero $<1$ rappresentato in binario come $.b_1b_2b_3\ldots$ \emph{(10/01/20)}
\item \qid{9.4} Enuncia il numero di bit emessi da Huffman e da Arithmetic dato un testo $T$ di
lunghezza $n$ ed entropia empirica $H$; commenta quando conviene comprimere con l'uno o con l'altro,
tenendo conto anche dell'efficienza in tempo. \emph{(06/03/25, 03/06/25)}
\item \qid{9.5} Dimostra che due testi $T_1$, $T_2$, l'uno permutazione dell'altro, sono compressi
con lo stesso numero di bit dalla codifica aritmetica. \emph{(23/06/25)}
\item \qid{9.6} Perch\'e la codifica aritmetica non emette l'estremo sinistro dell'intervallo finale,
ma parte dall'elemento centrale e lo tronca? \emph{(07/09/21)}
\end{itemize}

\section{Strutture dati compresse}

\begin{itemize}[leftmargin=1.4em]
\item \qid{10.1}\hot Descrivi la struttura dati per supportare $\texttt{Rank}$ su un array binario
$B[1,n]$ in tempo $\OO{1}$ e valutane lo spazio in bit. \emph{(16/01/23, 05/06/23, 02/11/23)}
\item \qid{10.2}\hot Dimostra la complessit\`a in tempo e in spazio (in bit) della soluzione che
implementa $\texttt{Rank}$ in $\OO{1}$ e $o(m)$ bit su $B[1,m]$. \emph{(12/12/23, 10/01/25, 04/02/25)}
\item \qid{10.3} (Riferimento) Codifica succinta di alberi binari in $B[1,2n+1]$ bit e navigazione
con Rank/Select; LOUDS per alberi generali. \emph{(base degli esercizi: 11/11/20, 15/12/21,
25/07/22, 16/01/23, 10/02/23, 24/07/23, 12/12/23, 09/02/24, 10/01/25, 04/02/25, 06/03/25, 03/06/25, 23/06/25)}
\end{itemize}

\section{Domande di progettazione}

\begin{itemize}[leftmargin=1.4em]
\item \qid{11.1} \`E data una sequenza $S$ di $n$ coppie $\langle$stringa, occ$\rangle$, dove pi\`u
coppie possono riferirsi alla stessa stringa e \emph{occ} \`e un valore satellite intero; sia $s$ il
numero di stringhe distinte. Progetta un algoritmo I/O-efficiente che calcola la stringa con massima
somma dei valori satellite, valutandone la complessit\`a I/O nei tre casi: (a) $M$ molto pi\`u
piccola di $s$; (b) $M=\OO{s\log n}$ bit ammettendo errori in probabilit\`a; (c) $M$ pu\`o contenere
$\OO{s}$ puntatori a stringhe e $\OO{s}$ caratteri e contatori (hint: soluzione trie-based, ma non un
trie classico). \emph{(07/02/20)}
\item \qid{11.2} Dato un testo $T[1,n]$ indicizzato da un Suffix Array, progetta un algoritmo che,
dati $P[1,p]$ e due interi positivi $k$ e $q$, stabilisce se esiste un range di $k$ posizioni
contigue $T[i,i+k-1]$ in cui \emph{iniziano} almeno $q$ occorrenze di $P$; stampa $i$ se esiste,
altrimenti $-1$. \emph{(02/02/22)}
\item \qid{11.3} \`E dato un insieme di coppie $\langle$frequenza, parola$\rangle$. Progetta una
struttura dati e un algoritmo che, dati un range di frequenze $[f_{\min},f_{\max}]$ e una parola
$W$, restituiscono tutte le parole lessicograficamente maggiori di $W$ con frequenza nel range;
discuti complessit\`a in tempo e spazio. \emph{(11/11/24)}
\end{itemize}


\newpage
\part{Risoluzioni dettagliate}
\section{Modelli e complessit\`a I/O --- risoluzioni}

\subsection*{\qid{1.1}\hot\ MergeSort multi-way: complessit\`a I/O}

\begin{domanda}
Enuncia la complessit\`a I/O del MergeSort multi-way in funzione di $M$, $B$, $N$.
\emph{(09/02/24, 27/05/24, 04/02/25)}
\end{domanda}

\textbf{Enunciato.} Con memoria interna $M$, pagina di disco $B$ e $N$ item da ordinare:
\[
\boxed{\ \text{MergeSort multi-way} = \OO{\frac{N}{B}\,\log_{M/B}\frac{N}{M}} \text{ I/O}\ }
\]
che coincide con il bound ottimo $sort(N)=\TT{\frac{N}{B}\log_{M/B}\frac{N}{M}}$.

\textbf{Da dove viene la formula (serve saperla giustificare).} L'algoritmo ha due fasi.

\emph{Fase 1 --- formazione dei run.} Si legge un blocco di $M$ item alla volta, lo si ordina in
memoria interna (con qualunque algoritmo ottimo, oppure con QuickSort/HeapSort) e lo si riscrive su
disco come \emph{run} ordinato. Ogni item \`e letto e scritto una volta: $\OO{N/B}$ I/O. Alla fine
ci sono $N/M$ run ordinati di lunghezza $M$.

\emph{Fase 2 --- merge a $k$ vie.} La memoria pu\`o tenere $M/B$ pagine. Si fondono
\kw{$k=\frac{M}{B}-1$ run alla volta}: una pagina di buffer per ciascuno dei $k$ run in ingresso
pi\`u una pagina di output (quando un buffer di input si svuota si carica la pagina successiva del
suo run; quando l'output si riempie si scarica su disco). Il minimo corrente fra i $k$ run si trova
con un min-heap di $k$ elementi ($\OO{\log k}$ tempo per estrazione, ma qui contiamo solo gli I/O).
\kw{Ogni passata di merge legge e scrive tutti gli $N$ item} ($\OO{N/B}$ I/O) e riduce il numero di
run di un fattore $k=\Theta(M/B)$. Il numero di passate per passare da $N/M$ run a 1 \`e dunque
\[
\log_k\frac{N}{M}=\log_{M/B}\frac{N}{M}.
\]
Totale: $\OO{\frac{N}{B}}$ (fase 1) $+\;\OO{\frac{N}{B}\log_{M/B}\frac{N}{M}}$ (fase 2), cio\`e la
formula enunciata. \kw{Il punto qualificante \`e la base del logaritmo: $M/B$ invece di 2} (del
MergeSort binario). Con valori realistici il fattore $\log_{M/B}(N/M)$ vale 1 o 2: in pratica
``due passate sul disco''.

\subsection*{\qid{1.2} Upper bound: sorting, permuting, striping}

\begin{domanda}
Enuncia gli upper bound per: ordinare $n$ item atomici in un modello a due livelli con memoria
interna $M$, pagina $B$ e 1 disco; permutarli nello stesso modello con 1 disco; ordinarli con $D$
dischi usando il disk striping. \emph{(10/02/23)}
\end{domanda}

\textbf{Sorting, 1 disco.} $\OO{\frac{n}{B}\log_{M/B}\frac{n}{M}}$ I/O (MergeSort multi-way,
\qid{1.1}).

\textbf{Permuting, 1 disco.}
\[
\boxed{\ \OO{\min\{\,n,\; sort(n)\,\}} \text{ I/O}\ }
\]
Le due alternative: (i) muovere gli item uno alla volta nella posizione finale, $\OO{1}$ I/O
ciascuno, totale $\OO{n}$; (ii) attaccare a ogni item la sua posizione di destinazione come chiave
e ordinare con MergeSort multi-way, totale $\OO{sort(n)}$. Si sceglie il minimo. (Il lower bound
corrispondente \`e $\Om{\min\{n,\frac{n}{B}\log_{M/B}\frac{n}{B}\}}$, e per i valori realistici dei
parametri i due problemi costano uguale: si veda \qid{1.5}.)

\textbf{Sorting, $D$ dischi con disk striping.} $\OO{\frac{n}{DB}\log_{M/(DB)}\frac{n}{M}}$ I/O:
si veda \qid{1.3}.

\subsection*{\qid{1.3}\hot\ Disk striping e MergeSort su $D$ dischi (con prova)}

\begin{domanda}
Descrivi la tecnica del disk striping ed enuncia il bound I/O del MergeSort quando lo si usa per
ordinare $n$ item in una memoria a due livelli con $D$ dischi, memoria $M$, pagina $B$.
\emph{(07/02/20 con prova, 09/02/24, 04/02/25, 03/06/25)}
\end{domanda}

\textbf{La tecnica.} Il disk striping fa lavorare $D$ dischi come \kw{un unico disco logico con
pagine di taglia $D\cdot B$}. La pagina logica $i$ \`e composta dalle $D$ pagine fisiche di indice
$i$, una per disco (la ``striscia'' $i$-esima): i primi $B$ item sulla pagina $i$ del disco 1, i
successivi $B$ sulla pagina $i$ del disco 2, ecc. Poich\'e i $D$ dischi possono leggere/scrivere in
parallelo, un singolo I/O logico trasferisce $D\cdot B$ item al costo (in tempo) di un I/O fisico.
Il vantaggio \`e la semplicit\`a: qualunque algoritmo pensato per 1 disco funziona immutato.

\textbf{Bound (prova).} Si esegue il MergeSort multi-way di \qid{1.1} sul disco logico, cio\`e
\kw{sostituendo $B\to DB$} nell'analisi. Contando gli I/O logici:
\[
\boxed{\ \OO{\frac{n}{DB}\,\log_{M/(DB)}\frac{n}{M}} \text{ I/O}\ }
\]
Ogni I/O logico corrisponde a $D$ I/O fisici \emph{simultanei}, quindi questa \`e anche la misura
del tempo di I/O con $D$ dischi in parallelo.

\subsection*{\qid{1.4} Confronto con il bound ottimo su $D$ dischi}

\begin{domanda}
Enuncia la complessit\`a I/O di ordinare $n$ item con memoria $M$, pagina $B$ e $D$ dischi;
confrontala con il bound ottimo e discuti le differenze. \emph{(07/02/20, 10/01/25)}
\end{domanda}

Il bound ottimo per ordinare con $D$ dischi \`e
\[
\boxed{\ \TT{\frac{n}{DB}\,\log_{M/B}\frac{n}{DB}}\ \text{I/O}\ }
\]
raggiunto da algoritmi pi\`u sofisticati (es.\ Greed Sort), che tengono la base del logaritmo a
$M/B$ pur dividendo il volume di dati fra i dischi.

\textbf{Differenze.} Il disk striping (\qid{1.3}) ottiene lo stesso fattore $\frac{n}{DB}$ davanti
al logaritmo, ma \kw{la base del logaritmo peggiora da $M/B$ a $M/(DB)$}: usando pagine logiche di
taglia $DB$, il fan-out del merge scende da $\Theta(M/B)$ a $\Theta(M/(DB))$, cio\`e si fondono
meno run per passata e servono pi\`u passate. Il rapporto fra i due bound \`e
\[
\frac{\log_{M/(DB)}(n/M)}{\log_{M/B}(n/DB)}\;\approx\;\frac{\log(M/B)}{\log(M/(DB))},
\]
che vale $\approx 1$ quando $DB\ll M$ (\kw{pochi dischi: striping quasi ottimo}) e diverge quando
$DB\to M$ (con $D$ grande il merge diventa quasi binario: \kw{striping non ottimale per $D$
grande}). In pratica si usa comunque lo striping per la sua semplicit\`a.

\subsection*{\qid{1.5} Sorting vs permuting: RAM vs modello a due livelli}

\begin{domanda}
Dati $n$ item atomici, scrivi la complessit\`a in tempo dei problemi di sorting e di permuting nel
modello RAM, poi la loro complessit\`a I/O nel modello a due livelli (pagina $B$, memoria $M$);
commenta se i due problemi sono computazionalmente equivalenti in senso asintotico e per quali
valori dei parametri. \emph{(orale 25/07/22)}
\end{domanda}

\textbf{Modello RAM.} Sorting (comparison-based): $\TT{n\log n}$. Permuting: $\TT{n}$ (si muove
ogni item direttamente in posizione finale). \kw{In RAM permutare \`e strettamente pi\`u facile}
di un fattore $\log n$.

\textbf{Modello a due livelli.} Sorting: $\TT{\frac{n}{B}\log_{M/B}\frac{n}{M}}$ I/O.
Permuting: $\TT{\min\{n,\;\frac{n}{B}\log_{M/B}\frac{n}{B}\}}$ I/O.

\textbf{Quando coincidono.} Il permuting costa quanto il sorting a meno che il termine ``naive''
$n$ non sia il minimo, cio\`e a meno che $B<\log_{M/B}\frac{n}{B}$. Con i valori reali dei
parametri ($B$ dell'ordine di $10^3$--$10^4$ item, mentre $\log_{M/B}(n/B)$ vale pochissime unit\`a
anche per dataset enormi) risulta sempre $B\gg\log_{M/B}(n/B)$: quindi \kw{nel modello a due
livelli sorting e permuting costano asintoticamente uguale}. \`E la differenza concettuale pi\`u
importante col modello RAM: su disco non si riesce a sfruttare la conoscenza della permutazione
finale, perch\'e muovere item singoli paga comunque un I/O intero a spostamento; tanto vale
ordinare. Morale del corso: \kw{in memoria esterna anche i problemi ``facili'' costano come il
sorting}, e $sort(n)$ \`e l'unit\`a di misura di riferimento.

% ==================================================================
\section{Ordinamento --- risoluzioni}

\subsection*{\qid{2.1}\hot\ Snow Plow: lunghezza attesa del run \`e $2M$}

\begin{domanda}
Dimostra che la lunghezza attesa della sequenza ordinata (run) prodotta da Snow Plow \`e $2M$.
Varianti chieste: quanto vale se la probabilit\`a di finire nell'unsorted bucket \`e $\tfrac14$
invece di $\tfrac12$? (15/01/24) Se \`e $\tfrac18$? (21/06/24) Se la probabilit\`a di andare nello
heap \`e $\tfrac34$? (11/11/20)
\end{domanda}

\textbf{L'algoritmo (serve descriverlo prima della prova).} Snow Plow genera i run ordinati della
fase 1 del MergeSort usando una memoria di $M$ celle divisa fra un min-heap $\mathcal{H}$ e un
buffer non ordinato $\mathcal{U}$. Una \emph{fase} (= la produzione di un run) funziona cos\`i:
all'inizio $\mathcal{H}$ contiene $M$ item e $\mathcal{U}$ \`e vuoto. Ripetutamente: si estrae il
minimo $m$ da $\mathcal{H}$ e lo si scrive nel run di output; si legge il prossimo item $x$ dallo
stream: se $x\ge m$ pu\`o ancora far parte del run corrente e va in $\mathcal{H}$, altrimenti \`e
``troppo piccolo'' (uscirebbe fuori ordine) e va in $\mathcal{U}$. La memoria resta sempre piena:
$|\mathcal{H}|+|\mathcal{U}|=M$. La fase termina quando $\mathcal{H}$ si svuota (equivalentemente
$\mathcal{U}$ si \`e riempito con $M$ item): i contenuti di $\mathcal{U}$ diventano lo heap della
fase successiva. Ogni run prodotto ha lunghezza $\ge M$.

\textbf{Teorema.} Se ogni item letto va in $\mathcal{U}$ con probabilit\`a $p=\tfrac12$, la
lunghezza attesa del run \`e $2M$.

\textbf{Prova.} Sia $\tau$ la lunghezza del run prodotto in una fase. \kw{Ogni item scritto in
output provoca la lettura di un item dallo stream}, quindi durante la fase vengono letti
esattamente $\tau$ item. Il run \`e formato da: gli $M$ item presenti nello heap all'inizio della
fase, pi\`u tutti gli item letti durante la fase che sono entrati in $\mathcal{H}$ (ciascuno di
essi, essendo $\ge$ del minimo corrente, verr\`a prima o poi estratto e scritto nello stesso run).
In valore atteso gli item letti che entrano nello heap sono $(1-p)\tau$. Dunque
\[
\tau \;=\; M+(1-p)\,\tau \qquad\Longrightarrow\qquad p\,\tau=M
\qquad\Longrightarrow\qquad \boxed{\;\tau=\frac{M}{p}\;}
\]
(passando ai valori attesi e usando la linearit\`a). Con $p=\tfrac12$: $\tau=2M$. $\blacksquare$

\textbf{Varianti d'esame} (la formula $\tau=M/p$ le risolve tutte):
\begin{itemize}[leftmargin=1.6em,itemsep=1pt]
\item $p=\tfrac14$ (15/01/24): \kw{$\tau=4M$};
\item $p=\tfrac18$ (21/06/24): \kw{$\tau=8M$};
\item probabilit\`a di andare nello \emph{heap} $=\tfrac34$ (11/11/20): allora
$p=1-\tfrac34=\tfrac14$ e \kw{$\tau=4M$}. \kw{Attenzione}: qui $p$ \`e la probabilit\`a
dell'unsorted bucket; se il testo d\`a quella dello heap, va complementata.
\end{itemize}
\textbf{Intuizione.} Pi\`u \`e raro che un item sia ``troppo piccolo'' (piccolo $p$), pi\`u a lungo
lo heap resta rifornito e pi\`u lungo \`e il run. Casi limite: input gi\`a ordinato $\Rightarrow$
$p=0$ $\Rightarrow$ un unico run; input ordinato al contrario $\Rightarrow p=1$ $\Rightarrow$ run
di lunghezza esattamente $M$. Effetto sul MergeSort: run attesi $n/2M$, cio\`e
$\OO{\frac{n}{B}\log_{M/B}\frac{n}{2M}}$ I/O in media.

\subsection*{\qid{2.2}\hot\ Lower bound per l'ordinamento di stringhe}

\begin{domanda}
Enuncia e dimostra il lower bound (comparison-based) al tempo per ordinare $n$ stringhe di
lunghezza totale $N$ su alfabeto di taglia $\sigma$. \emph{(10/02/23, 05/06/23, orale 25/07/22)}
\end{domanda}

\textbf{Notazione.} Per ogni stringa $s$ del dataset $S$ sia $d_s$ la lunghezza del suo
\emph{prefisso distinguente}, cio\`e il pi\`u corto prefisso di $s$ che non \`e prefisso di
nessun'altra stringa di $S$; sia $d=\sum_{s\in S} d_s$ ($d\le N$).

\textbf{Teorema.} Ogni algoritmo comparison-based che ordina $n$ stringhe distinte richiede
\[
\boxed{\ \Om{d+n\log n}\ }
\]
confronti di singoli caratteri, nel caso peggiore.

\textbf{Prova.} Si dimostrano separatamente i due termini; il massimo dei due \`e un lower bound,
e poich\'e $\max\{a,b\}\ge\frac{a+b}{2}$ lo \`e anche la somma (a meno di un fattore 2).

\emph{(i) $\Om{n\log n}$.} Argomento dell'albero di decisione: l'algoritmo deve distinguere gli
$n!$ ordinamenti possibili delle stringhe; ogni confronto fra caratteri ha esito in un insieme di
taglia costante (es.\ $<,=,>$), quindi l'albero delle computazioni ha grado costante e servono
$\Om{\log(n!)}=\Om{n\log n}$ confronti su qualche ramo.

\emph{(ii) $\Om{d}$.} \kw{L'algoritmo deve leggere almeno tutti i caratteri dei prefissi
distinguenti.} Supponiamo per assurdo che su un input $S$ l'algoritmo termini senza aver letto il
carattere $s[i]$ con $i\le d_s$ per qualche $s$. Poich\'e $i\le d_s$, il prefisso $s[1,i]$ non \`e
ancora distinguente: esiste un input alternativo, indistinguibile per l'algoritmo (coincide con $S$
su tutti i caratteri letti), in cui il carattere in posizione $i$ cambia l'ordine relativo di $s$.
L'algoritmo produce lo stesso output su due input con ordinamenti diversi: contraddizione. Quindi
$\Om{d}$ letture. $\blacksquare$

\textbf{Commento.} $n\log n$ \`e il costo ``informativo'' dell'ordinamento; $d$ \`e il costo di
``vedere'' abbastanza caratteri. \kw{Non compare $N$}: i caratteri dopo il prefisso distinguente
sono irrilevanti. Il Multikey QuickSort (\qid{2.3}) raggiunge questo bound: il bound \`e stretto.

\subsection*{\qid{2.3}\hot\ Multikey QuickSort: pseudocodice, complessit\`a, ottimalit\`a}

\begin{domanda}
Scrivi lo pseudocodice del Multikey QuickSort su $n$ stringhe di lunghezza totale $N$; enuncia e
dimostra la sua complessit\`a temporale e mostra che \`e time-optimal rispetto al lower bound
\qid{2.2}. \emph{(10/02/23, 05/06/23, 07/09/23, orale 25/07/22)}
\end{domanda}

\textbf{Idea.} \`E il QuickSort adattato alle stringhe: si partiziona rispetto a \kw{un solo
carattere} (quello in posizione $i$, dove $i$ \`e la lunghezza del prefisso comune gi\`a
accertato) in tre gruppi $<,=,>$, e \kw{solo nel gruppo $=$ si avanza al carattere successivo}.

\begin{verbatim}
MultikeyQS(R, i):                    // R insieme di stringhe, i posizione corrente
  if |R| <= 1: return R
  scegli una stringa pivot p da R uniformemente a caso;  c = p[i]
  R<  = { s in R : s[i] < c }
  R=  = { s in R : s[i] = c }
  R>  = { s in R : s[i] > c }
  return MultikeyQS(R<, i) ++ MultikeyQS(R=, i+1) ++ MultikeyQS(R>, i)
\end{verbatim}

Invariante: tutte le stringhe di $R$ condividono il prefisso $p[1,i-1]$, quindi confrontare il
solo carattere $i$ \`e corretto; la concatenazione $R_<\,R_=\,R_>$ d\`a l'ordine lessicografico.

\textbf{Teorema.} Con pivot casuale, il tempo atteso \`e
\[
\boxed{\ \OO{d+n\log n}\ \text{atteso}\ }
\]
dove $d$ \`e la somma dei prefissi distinguenti.

\textbf{Prova.} Contiamo i confronti di carattere addebitandoli alla stringa $s$ che li subisce.
Ogni confronto colloca $s$ in $R_<$, $R_=$ o $R_>$:
\begin{itemize}[leftmargin=1.6em,itemsep=1pt]
\item \emph{Confronti di tipo ``$=$''}: $s$ finisce in $R_=$ e l'indice $i$ avanza di 1. Questo
pu\`o accadere al pi\`u $d_s+\OO{1}$ volte: \kw{appena $i$ supera $d_s$ il prefisso di $s$ \`e
unico}, quindi $s$ resta da sola nel suo gruppo e la sua ricorsione termina. Totale: $\OO{d}$.
\item \emph{Confronti di tipo ``$<$'' o ``$>$''}: $s$ finisce in $R_<$ o $R_>$ e $i$ non avanza.
Questi confronti sono, in valore atteso, $\OO{\log n}$ per stringa: \`e l'analisi del QuickSort
classico ristretta al carattere $i$-esimo. Un pivot casuale \`e, con probabilit\`a $\ge\tfrac12$,
``centrale'' (n\'e fra i pi\`u piccoli $\tfrac14$ n\'e fra i pi\`u grandi $\tfrac14$ dei caratteri
$i$-esimi del gruppo): in tal caso $|R_<|,|R_>|\le\tfrac34|R|$. Quindi a ogni confronto di questo
tipo, con probabilit\`a costante il gruppo di $s$ si riduce di un fattore costante: pu\`o accadere
$\OO{\log n}$ volte. Totale: $\OO{n\log n}$.
\end{itemize}
Sommando: $\OO{d+n\log n}$ atteso. $\blacksquare$

\textbf{Ottimalit\`a.} \kw{Coincide col lower bound \qid{2.2}}: il Multikey QuickSort \`e
time-optimal (in attesa). Vantaggi ingegneristici: non spezza le stringhe (localit\`a), partizione
ternaria in place, non richiede di conoscere l'alfabeto.

\subsection*{\qid{2.4}\hot\ LSD RadixSort: complessit\`a e correttezza}

\begin{domanda}
Abbozza l'algoritmo LSD RadixSort per $n$ chiavi di $b$ bit ciascuna; enuncia e dimostra la sua
complessit\`a temporale; dimostra la correttezza dell'ordinamento prodotto (caso stringhe
binarie). Variante: complessit\`a del RadixSort ottimo su $n$ stringhe binarie di lunghezza totale
$N$. \emph{(16/01/23, 05/06/23, 11/07/24, 10/01/25, 23/06/25)}
\end{domanda}

\textbf{Algoritmo.} Le $n$ chiavi di $b$ bit sono viste come sequenze di $b/r$ \emph{cifre} di $r$
bit (cifre in $[0,2^r)$). L'algoritmo fa $b/r$ passate, \kw{dalla cifra meno significativa alla
pi\`u significativa}; ogni passata riordina l'intera sequenza rispetto alla cifra corrente con
CountingSort, che \`e \kw{stabile}:

\begin{verbatim}
LSD-RadixSort(S, b, r):
  for j = 1 .. b/r:                      // j = 1 e' la cifra MENO significativa
      S = CountingSortStabile(S, cifra j)   // O(n + 2^r) tempo
  return S
\end{verbatim}

\textbf{Complessit\`a.} Ogni passata costa $\OO{n+2^r}$ (istogramma delle $2^r$ cifre + somme
prefisse + ricollocamento). Totale:
\[
\boxed{\ \OO{\frac{b}{r}\,(n+2^r)}
\;\xrightarrow{\;r=\log_2 n\;}\;
\OO{\frac{b}{\log n}\cdot n}\ }
\]
\kw{La scelta ottima \`e $r=\log_2 n$} (cos\`i $2^r=n$ non domina): per chiavi di $b=\OO{\log n}$
bit il costo \`e $\OO{n}$ --- batte il lower bound comparison-based perch\'e \kw{il RadixSort non
\`e comparison-based}. Spazio: $\OO{n+2^r}$ ausiliario. \emph{Variante} (05/06/23), $n$ stringhe
binarie di lunghezza totale $N$ (cio\`e $b=N/n$): tempo
$\OO{\frac{nb}{\log n}}=\OO{\frac{N}{\log n}}$ (per alfabeto $\sigma$:
$\OO{\frac{N\log\sigma}{\log n}}$).

\textbf{Correttezza (induzione + stabilit\`a).} Invariante: \kw{dopo la passata sulla cifra $j$,
le chiavi sono ordinate rispetto alle ultime $j$ cifre} (le $j$ meno significative).

\emph{Base} ($j=0$): banale.

\emph{Passo.} Vale l'invariante per $j-1$; eseguiamo la passata $j$ e consideriamo due chiavi
$\alpha,\beta$ qualunque:
\begin{itemize}[leftmargin=1.6em,itemsep=1pt]
\item \emph{cifra $j$ diversa}: il CountingSort le ordina secondo la cifra $j$, che \`e la pi\`u
significativa fra le ultime $j$: ordine corretto qualunque fosse prima;
\item \emph{cifra $j$ uguale}: il loro ordine \`e deciso dalle ultime $j-1$ cifre; per ipotesi
induttiva erano gi\`a nell'ordine giusto, e \kw{la stabilit\`a del CountingSort preserva tale
ordine}.
\end{itemize}
Dopo la passata $j=b/r$: ordinate rispetto a tutte le cifre. $\blacksquare$

\kw{Nota: la stabilit\`a \`e essenziale} --- con un ordinamento interno non stabile l'invariante si
rompe sui pareggi della cifra corrente.

\subsection*{\qid{2.5}\hot\ Bounded QuickSort: stack $\OO{\log n}$}

\begin{domanda}
Scrivi e commenta lo pseudocodice del Bounded QuickSort, cio\`e la variante che garantisce uno
stack di taglia $\OO{\log n}$ ($\OO{\log n}$ chiamate ricorsive annidate per qualunque input).
\emph{(10/01/25, 04/02/25, 23/06/25)}
\end{domanda}

\textbf{Problema.} Il QuickSort ricorsivo standard, nel caso peggiore (partizioni sbilanciate), ha
profondit\`a di ricorsione $\OO{n}$: rischio di stack overflow. La variante \emph{bounded} elimina
una delle due chiamate ricorsive con l'iterazione (tail recursion removal), \kw{ricorrendo sempre
e solo sul lato pi\`u piccolo}:

\begin{verbatim}
BoundedQuickSort(S, i, j):          // ordina S[i..j]
  while i < j:
      p = Partition(S, i, j)        // pivot (casuale o mediana di campione)
      if (p - i) < (j - p):         // lato sinistro piu' piccolo
           BoundedQuickSort(S, i, p-1)
           i = p + 1                // il lato destro si tratta iterativamente
      else:                         // lato destro piu' piccolo (o uguale)
           BoundedQuickSort(S, p+1, j)
           j = p - 1                // il lato sinistro si tratta iterativamente
\end{verbatim}

\textbf{Commento / prova del bound sullo stack.} La chiamata ricorsiva \`e sempre eseguita sul
lato di taglia $\le\frac{j-i+1}{2}$. Quindi, \kw{scendendo di un livello di ricorsione la taglia
del sottoproblema almeno si dimezza}: dopo $t$ livelli annidati \`e $\le n/2^t$, perci\`o
\[
\boxed{\ t\le\log_2 n \ \text{ livelli di stack, worst case, per qualunque input}\ }
\]
Il lato grande non consuma stack: \`e gestito dal ciclo \texttt{while} restringendo $[i,j]$. Il
tempo resta quello del QuickSort ($\OO{n\log n}$ atteso con pivot casuale, $\OO{n^2}$ worst case):
\kw{il bound riguarda lo spazio, non il tempo}.

% ==================================================================
\section{Random Sampling --- risoluzioni}

\subsection*{\qid{3.1}--\qid{3.2}\hot\ Reservoir Sampling: pseudocodice e correttezza}

\begin{domanda}
\qid{3.1} Scrivi (solo) lo pseudocodice del Reservoir Sampling. \emph{(06/03/25, 03/06/25)}\\
\qid{3.2} Data una sequenza $S$ di $n$ item ($n$ ignoto) e un intero $m$ (eventualmente $m<n/2$):
mostra lo pseudocodice del Reservoir Sampling e dimostrane la correttezza, cio\`e che ogni item
\`e campionato con probabilit\`a $m/n$. \emph{(02/11/23, 11/07/24)}
\end{domanda}

\textbf{Pseudocodice.}
\begin{verbatim}
ReservoirSampling(S, m):           // n = |S| ignoto a priori
  R[1..m] = primi m elementi di S
  i = m
  while (S ha un altro elemento x):
      i = i + 1
      j = intero uniforme in [1, i]
      if j <= m:  R[j] = x         // x entra, R[j] esce
  return R
\end{verbatim}

Una sola passata, tempo $\OO{n}$, spazio $\OO{m}$; in ogni istante $R$ \`e un campione uniforme
del prefisso letto: adatto a stream di lunghezza sconosciuta.

\textbf{Teorema (correttezza).} Al termine, ogni elemento di $S$ appartiene a $R$ con
probabilit\`a esattamente $m/n$.

\textbf{Prova.} Consideriamo l'elemento in posizione $t$.

\emph{Caso $t>m$.} L'elemento \kw{entra} in $R$ al passo $t$ con probabilit\`a
$\Pr[j\le m]=\frac{m}{t}$. Per ogni passo successivo $i>t$, l'elemento in $R$ viene \kw{espulso}
sse l'elemento $i$-esimo entra (prob.\ $\frac{m}{i}$) \emph{e} sceglie proprio la sua cella
(prob.\ $\frac{1}{m}$): probabilit\`a $\frac{m}{i}\cdot\frac{1}{m}=\frac{1}{i}$; quindi
\kw{sopravvive al passo $i$ con probabilit\`a $\frac{i-1}{i}$}. Per indipendenza delle estrazioni:
\[
\Pr[t\in R] \;=\; \frac{m}{t}\cdot\!\!\prod_{i=t+1}^{n}\frac{i-1}{i}
\;=\;\frac{m}{t}\cdot\frac{t}{t+1}\cdot\frac{t+1}{t+2}\cdots\frac{n-1}{n}
\;=\;\frac{m}{t}\cdot\frac{t}{n}\;=\;\boxed{\ \frac{m}{n}\ }
\]
per \kw{prodotto telescopico}.

\emph{Caso $t\le m$.} Entra con probabilit\`a 1 e deve sopravvivere ai passi $i=m+1,\dots,n$:
$\Pr[t\in R]=\prod_{i=m+1}^{n}\frac{i-1}{i}=\frac{m}{n}$.

In entrambi i casi $\frac{m}{n}$: campionamento uniforme. $\blacksquare$

\subsection*{\qid{3.3} Sampling con $n$ noto (selection sampling) e correttezza per $m=1$}

\begin{domanda}
Data una sequenza $S$ di $n$ item con $n$ noto e un intero $m$, scrivi lo pseudocodice
dell'algoritmo di sampling che estrae $m$ item uniformemente a caso, e dimostrane la correttezza
(probabilit\`a $m/n$) nel caso $m=1$. \emph{(06/11/23; simulazioni in 13/06/22, 10/02/23,
07/09/23)}
\end{domanda}

Quando $n$ \`e noto si decide \emph{al volo}, per ogni elemento, se selezionarlo, con probabilit\`a
che dipende da quanti ne mancano:

\begin{verbatim}
KnownN-Sampling(S, n, m):
  s = 0                              // gia' selezionati
  for i = 1 .. n:
      resta  = n - i + 1             // elementi ancora da esaminare (incluso i)
      serve  = m - s                 // elementi ancora da selezionare
      con probabilita' serve/resta:  // es.: estrai p uniforme in [0,1],
           seleziona S[i];  s = s+1  //       seleziona se p <= serve/resta
      if s == m: break
  return i selezionati
\end{verbatim}

\kw{Non pu\`o selezionare pi\`u di $m$ elementi} (con $s=m$ la probabilit\`a \`e 0) \kw{n\'e meno}
(quando $\textit{serve}=\textit{resta}$ la probabilit\`a \`e 1 per tutti i rimanenti). Una
passata, $\OO{n}$ tempo, $\OO{1}$ spazio. \emph{Nota per le simulazioni d'esame}: i valori $p_i$
dati nel testo si confrontano con la soglia corrente $\textit{serve}/\textit{resta}$
(si seleziona quando $p_i\le\textit{serve}/\textit{resta}$).

\textbf{Correttezza per $m=1$.} L'elemento $t$ \`e selezionato sse nessuno dei precedenti lo \`e
stato e la moneta del passo $t$ ha successo (prob.\ $\frac{1}{n-t+1}$):
\[
\Pr[t \text{ selezionato}]
=\prod_{i=1}^{t-1}\Bigl(1-\frac{1}{n-i+1}\Bigr)\cdot\frac{1}{n-t+1}
=\prod_{i=1}^{t-1}\frac{n-i}{n-i+1}\cdot\frac{1}{n-t+1}
=\frac{n-t+1}{n}\cdot\frac{1}{n-t+1}=\boxed{\ \frac{1}{n}\ }
\]
di nuovo per telescopia: uniforme. $\blacksquare$ (Il caso generale \`e analogo per induzione, e il
numero di selezionati \`e esattamente $m$ per costruzione.)

\textbf{Confronto col Reservoir.} Serve conoscere $n$, ma niente serbatoio n\'e sostituzioni: ogni
elemento \`e deciso definitivamente al suo passaggio.

% ==================================================================
\section{Set Intersection --- risoluzioni}

Setting comune: $L_1$ (detta anche $A$) di lunghezza $n$ e $L_2$ (detta $B$) di lunghezza $m\le n$,
ordinate crescenti. Riferimenti: merge completo $\OO{n+m}$ (buono per $m\approx n$); $m$ ricerche
binarie $\OO{m\log n}$ (buono per $m\ll n$). Gli algoritmi seguenti interpolano fra i due.

\subsection*{\qid{4.1}\hot\ Two-level intersection}

\begin{domanda}
Date due liste ordinate $L_1$, $L_2$ di lunghezze $n$ e $m$: descrivi l'algoritmo two-level (detto
anche two-level scan / two-level memory / two-level storage) per calcolarne l'intersezione;
enuncia e dimostra la sua complessit\`a. \emph{(09/02/24, 21/06/24, 11/11/24, 15/07/25)}
\end{domanda}

\textbf{Algoritmo.} Si fissa un parametro $L$ (taglia di blocco). Si divide $L_1$ in $h=n/L$
blocchi consecutivi di $L$ elementi e si costruisce la lista di primo livello
$L_1'=\{$primo elemento di ogni blocco$\}$, di lunghezza $n/L$.
\emph{Fase 1}: merge di $L_1'$ con $L_2$: costo $\OO{\frac{n}{L}+m}$; per ogni $b\in L_2$ questo
individua l'unico blocco di $L_1$ che pu\`o contenerlo. \emph{Fase 2}: per ogni blocco selezionato,
merge locale del blocco con i suoi candidati; \kw{ogni elemento di $L_2$ fa scandire al pi\`u un
blocco}: costo $\OO{mL+m}$ nel caso peggiore.

\textbf{Teorema.} Il costo \`e $\OO{\frac{n}{L}+mL+m}$; con la scelta ottima $L=\sqrt{n/m}$:
\[
\boxed{\ T(n,m)=\OO{m+\sqrt{n\,m}}\ }
\]

\textbf{Prova.} La somma delle due fasi \`e $f(L)=\frac{n}{L}+mL$ (pi\`u $\OO{m}$ additivo).
Minimizzo: $f'(L)=-\frac{n}{L^2}+m=0\iff L=\sqrt{n/m}$, dove $f(L)=2\sqrt{nm}$ (equivalentemente,
per AM--GM, $\frac{n}{L}+mL\ge2\sqrt{nm}$ con uguaglianza in $L=\sqrt{n/m}$). $\blacksquare$

\textbf{Commento.} $\sqrt{nm}$ \kw{batte sia il merge $\OO{n+m}$ (quando $m\ll n$) sia le ricerche
binarie $\OO{m\log n}$ (quando $m$ \`e grande)}. In I/O: $L_1$ \`e scandita per blocchi contigui,
quindi $\OO{\frac{n}{LB}+\frac{m}{B}+m}$ I/O con blocco $=$ pagina di disco (versione ``two-level
memory'' del 09/02/24: primi elementi in memoria, un fetch di blocco per elemento di $L_2$).

\subsection*{\qid{4.2} Mutual partition}

\begin{domanda}
Descrivi l'algoritmo mutual partition per l'intersezione di $L_1$, $L_2$, con la sua
complessit\`a. \emph{(11/07/24)}
\end{domanda}

\textbf{Algoritmo.} Si prende l'elemento mediano $b_{med}$ della lista \kw{corta} $L_2$; lo si
cerca con ricerca binaria in $L_1$, trovando la posizione $j$ tale che $L_1[1,j]<b_{med}\le
L_1[j+1,n]$ (se $b_{med}\in L_1$ lo si emette). Si ricorre sulle due met\`a:
$(L_1[1,j],\,L_2[1,med-1])$ e $(L_1[j+1,n],\,L_2[med+1,m])$. Base: lista vuota.

\textbf{Complessit\`a:} $\boxed{\OO{m\log\frac{2n}{m}}}$. Al livello $k$ della ricorsione ci sono
al pi\`u $2^k$ sottoproblemi; le porzioni di $L_1$ di uno stesso livello sono disgiunte e sommano
a $\le n$; ogni sottoproblema fa una ricerca binaria $\OO{1+\log n_{k,t}}$ sulla propria porzione.
\kw{Per concavit\`a del logaritmo (Jensen)} il costo del livello \`e
$\sum_{t}(1+\log n_{k,t})\le 2^k(1+\log\frac{n}{2^k})$. Sommando sui livelli $k=0,\dots,\log m$:
\[
\sum_{k=0}^{\log m}2^k\Bigl(1+\log\frac{n}{2^k}\Bigr)
=\OO{m}+\log n\sum_k 2^k-\sum_k k\,2^k
=\OO{m+m\log n-m\log m}=\OO{m\log\tfrac{2n}{m}},
\]
usando $\sum_{k\le K}2^k=\OO{m}$ e $\sum_{k\le K}k2^k=(K-1)2^{K+1}+2$ con $K=\log m$.
$\blacksquare$

\textbf{Commento.} \kw{\`E ottimo}: lower bound information-theoretic
$\Om{\log\binom{n}{m}}=\Om{m\log\frac{n}{m}}$ (bisogna distinguere gli interleaving possibili).
Svantaggio: accessi non sequenziali (ricorsione + binarie), localit\`a scarsa; il doubling
\qid{4.3} ottiene lo stesso bound con accesso iterativo.

\subsection*{\qid{4.3} Doubling / galloping search}

\begin{domanda}
Descrivi il doubling algorithm (binary search con salti esponenziali, ``galloping'') per
l'intersezione; enuncia e dimostra la sua complessit\`a. \emph{(15/01/24)}
\end{domanda}

\textbf{Algoritmo.} Si scandisce $L_2=b_1<\dots<b_m$ mantenendo in $L_1$ la posizione $i_{j-1}$
dell'ultimo attraversamento. Per il prossimo $b_j$: da $i_{j-1}$ si fanno \kw{salti esponenziali
$1,2,4,8,\dots$} finch\'e non si supera un elemento $\ge b_j$; detto $2^{k_j}$ l'ultimo salto, si
fa ricerca binaria di $b_j$ nell'intervallo di ampiezza $\Delta_j\le 2^{k_j}$ cos\`i individuato;
si emette $b_j$ se presente e si aggiorna $i_j$.

\textbf{Teorema.} Costo totale $\boxed{\OO{m\bigl(1+\log\frac{n}{m}\bigr)}}$.

\textbf{Prova.} Salti + binaria per $b_j$: $\OO{1+\log\Delta_j}$. Poich\'e i salti raddoppiano,
\kw{$\Delta_j\le 2\,(i_j-i_{j-1})$}: l'overshooting \`e al pi\`u il doppio del progresso reale,
quindi $\sum_j\Delta_j\le 2n$. \kw{Per concavit\`a del logaritmo (Jensen)}:
\[
\sum_{j=1}^{m}\OO{1+\log\Delta_j}
=\OO{m+m\log\Bigl(\tfrac{1}{m}\sum_j\Delta_j\Bigr)}
=\OO{m+m\log\tfrac{2n}{m}}. \;\blacksquare
\]

\textbf{Commento.} Stesso bound ottimo del mutual partition, ma con \kw{scansione monotona di
$L_1$}: localit\`a migliore (cache/prefetching), niente ricorsione. \`E l'algoritmo dei motori di
ricerca sulle posting list.

\subsection*{\qid{4.4} Intersezione di liste Elias--Fano con Access/GEQ}

\begin{domanda}
Scrivi lo pseudocodice per intersecare due liste ordinate $L_1$, $L_2$ (di $n_1$ e $n_2$ elementi)
codificate con Elias--Fano, usando come primitive $\texttt{Access}(L,i)$ e $\texttt{GEQ}(L,x)$;
valuta e commenta la complessit\`a. \emph{(14/01/22)}
\end{domanda}

Primitive ($\OO{1}$ con le strutture di rank/select sui bit alti, \qid{8.2}):
$\texttt{Access}(L,i)=$ $i$-esimo elemento; $\texttt{GEQ}(L,x)=$ posizione del pi\`u piccolo
elemento $\ge x$ (o $\bot$).

\begin{verbatim}
EF-Intersect(L1, L2):               // n1 = |L1|, n2 = |L2|
  i = 1
  x = Access(L1, 1)
  loop:
      j = GEQ(L2, x);   if j = NIL: return       // L2 esaurita
      y = Access(L2, j)
      if y == x:                                  // match
           output x
           i = i + 1;   if i > n1: return
           x = Access(L1, i)
      else:                                       // y > x: scavalca in L1
           i = GEQ(L1, y);   if i = NIL: return
           x = Access(L1, i)
\end{verbatim}

\textbf{Complessit\`a.} \kw{Ogni iterazione avanza strettamente almeno uno dei due puntatori} (su
match avanza $i$; altrimenti $x$ salta a un valore $\ge y>x$). Le iterazioni sono quindi
$\OO{n_1+n_2}$, e $\OO{\min(n_1,n_2)}$ chiamando GEQ dalla lista corta verso la lunga; ogni
operazione \`e $\OO{1}$: totale $\boxed{\OO{\min(n_1,n_2)}}$ \kw{senza mai decomprimere le liste}.
Vantaggio rispetto a \qid{4.3} su liste esplicite: spazio ($\approx\log\frac{u}{n}+2$ bit per
intero) e quindi meno I/O e cache miss; GEQ fa il ruolo dei salti esponenziali in $\OO{1}$.

\section{Hashing, Bloom Filter, Skip List, Treap --- risoluzioni}

\subsection*{\qid{5.1}\hot\ Hashing universale: definizione, esempio, prova}

\begin{domanda}
Definisci una classe di funzioni hash universali; fornisci un esempio e dimostrane
l'universalit\`a. \emph{(16/01/23, 07/09/23, 02/11/23, 10/01/25)}
\end{domanda}

\textbf{Definizione.} Una famiglia $\mathcal{H}\subseteq\{h:U\to[0,m)\}$ \`e \emph{universale} se
per ogni coppia di chiavi distinte $x\ne y$ di $U$:
\[
\boxed{\ \Pr_{h\in\mathcal{H}}\bigl[h(x)=h(y)\bigr]\;\le\;\frac{1}{m}\ }
\]
dove $h$ \`e estratta uniformemente da $\mathcal{H}$. Equivalentemente: il numero di funzioni
della famiglia che fanno collidere $x$ e $y$ \`e $\le|\mathcal{H}|/m$. \kw{La casualit\`a sta
nella scelta della funzione, non nelle chiavi}: la garanzia vale per ogni input, anche avversario.

\textbf{Esempio (prodotto scalare modulo primo).} Sia $m$ \kw{primo} e si scrivano le chiavi in
base $m$: $k=(k_0,k_1,\dots,k_{r-1})$ con $k_i\in[0,m)$. Per ogni vettore
$a=(a_0,\dots,a_{r-1})\in[0,m)^r$ si definisce
\[
h_a(k)\;=\;\Bigl(\sum_{i=0}^{r-1}a_i\,k_i\Bigr)\bmod m ,
\]
con $\mathcal{H}=\{h_a\}$, $|\mathcal{H}|=m^r$ (scegliere $h$ a caso $=$ scegliere $a$ a caso).

\textbf{Prova di universalit\`a.} Siano $x\ne y$; differiscono in almeno una cifra, diciamo la
$d$-esima: $x_d\ne y_d$. La collisione $h_a(x)=h_a(y)$ equivale a
\[
a_d\,(x_d-y_d)\;\equiv\;-\!\!\sum_{i\ne d}a_i\,(x_i-y_i) \pmod m .
\]
Poich\'e $m$ \`e primo e $x_d-y_d\not\equiv0$, \kw{$(x_d-y_d)$ \`e invertibile in $\mathbb{Z}_m$}:
fissati comunque gli altri $r-1$ coefficienti, esiste \kw{esattamente un} valore di $a_d$ che
soddisfa l'equazione:
\[
a_d\;\equiv\;-(x_d-y_d)^{-1}\!\!\sum_{i\ne d}a_i\,(x_i-y_i)\pmod m .
\]
Le funzioni che fanno collidere $x,y$ sono dunque $m^{r-1}$, e
$\Pr[h_a(x)=h_a(y)]=\frac{m^{r-1}}{m^{r}}=\frac{1}{m}$: universale. $\blacksquare$

\subsection*{\qid{5.2} Universale $+$ chaining: catene attese $\OO{1}$}

\begin{domanda}
Definisci una classe universale e mostra che, usata nell'hashing con chaining, ottiene lunghezza
attesa $\OO{1}$ per catena quando $m=\OO{n}$; da qui il calcolo del numero atteso di collisioni.
\emph{(06/11/23; esercizi collegati in 10/07/20 e 11/11/20)}
\end{domanda}

Definizione di classe universale: come \qid{5.1}. Sia $T[0,m-1]$ una tabella con chaining, $h$
estratta da una famiglia universale, $n$ chiavi inserite. Fissata una chiave $x$, sia $C_x$ il
numero di chiavi $y\ne x$ che collidono con $x$. \kw{Per linearit\`a dell'attesa e
universalit\`a}:
\[
\mathbb{E}[C_x]\;=\;\sum_{y\ne x}\Pr[h(y)=h(x)]\;\le\;\frac{n-1}{m}\;<\;\alpha:=\frac{n}{m}.
\]
Quindi il costo atteso di una ricerca \`e $\boxed{\OO{1+\alpha}}$, che \`e $\OO{1}$ non appena
$m=\Omega(n)$: \kw{$m=\OO{n}$ celle bastano per catene di lunghezza attesa costante}.
$\blacksquare$

\textbf{Esercizio collegato (10/07/20, 11/11/20): collisioni totali attese.} Il numero di coppie
collidenti \`e $C=\sum_{\{x,y\}}\mathbf{1}[h(x)=h(y)]$ su $\binom{n}{2}$ coppie:
\[
\mathbb{E}[C]\;\le\;\binom{n}{2}\frac{1}{m}\;=\;\frac{n(n-1)}{2m}.
\]
Con i numeri d'esame ($n=10$, range $[0,22]\Rightarrow m=23$):
$\mathbb{E}[C]\le\frac{10\cdot9}{2\cdot23}=\frac{45}{23}\approx1.96$.

\subsection*{\qid{5.3}\hot\ Cuckoo hashing: teorema dei cammini}

\begin{domanda}
Enuncia e dimostra il teorema sul cuckoo graph: per ogni coppia di posizioni $i,j$ della tabella
cuckoo, la probabilit\`a che il cammino minimo $i\to j$ abbia lunghezza $L$ \`e al pi\`u
$c^{-L}/r$, dove $r$ \`e il numero di posizioni della tabella, $n$ le chiavi memorizzate e
$r\ge 2cn$. Commenta l'impatto sul load factor. \emph{(10/01/20, 10/02/23, 03/06/25)}
\end{domanda}

\textbf{Contesto.} Nel cuckoo hashing ogni chiave $k$ ha due posizioni possibili
$h_1(k),h_2(k)$ in una tabella di $r$ celle; l'inserimento sfratta a catena. Il \emph{cuckoo
graph} ha come nodi le $r$ celle e, per ogni chiave, un arco fra le sue due posizioni. Le
prestazioni dipendono dall'assenza di cammini lunghi e cicli.

\textbf{Teorema.} Se le funzioni hash sono (idealmente) casuali e $r\ge 2cn$ per una costante
$c>1$, allora per ogni coppia di celle $i,j$ e ogni $L\ge1$:
\[
\boxed{\ \Pr\bigl[\exists\text{ cammino minimo }i\to j\text{ di lunghezza }L\bigr]
\;\le\;\frac{c^{-L}}{r}\ }
\]

\textbf{Prova (induzione su $L$).}

\emph{Base $L=1$.} Esiste un arco $(i,j)$ sse qualche chiave $k$ ha $\{h_1(k),h_2(k)\}=\{i,j\}$.
Per una singola chiave: $\Pr=\frac{2}{r^2}$ (due orientamenti). Per union bound sulle $n$ chiavi,
\kw{usando l'ipotesi $r\ge2cn$, cio\`e $n\le r/(2c)$}:
\[
\Pr[\exists\,\text{arco } i\text{--}j]\;\le\;\frac{2n}{r^{2}}
\;\le\;\frac{2}{r^{2}}\cdot\frac{r}{2c}\;=\;\frac{c^{-1}}{r}.
\]

\emph{Passo.} Un cammino minimo $i\to j$ di lunghezza $L$ \`e composto da un cammino minimo di
lunghezza $L-1$ da $i$ a una cella intermedia $\ell$, seguito da un arco $(\ell,j)$. Per union
bound su $\ell$, con ipotesi induttiva e indipendenza dell'ultimo arco:
\[
\Pr[\text{cammino }i\to j\text{ di lung. }L]
\;\le\;\sum_{\ell=1}^{r}\frac{c^{-(L-1)}}{r}\cdot\frac{c^{-1}}{r}
\;=\;r\cdot\frac{c^{-(L-1)}}{r}\cdot\frac{1}{c\,r}
\;=\;\frac{c^{-L}}{r}. \;\blacksquare
\]

\textbf{Conseguenze e load factor.} Probabilit\`a che $i$ e $j$ siano connessi (lunghezza
qualunque): $\sum_{L\ge1}\frac{c^{-L}}{r}=\frac{1}{r(c-1)}$. Da qui:
(i) \kw{insert atteso $\OO{1}$} (cammini attesi di lunghezza costante) e \kw{ricerca/delete
$\OO{1}$ worst case} (si guardano 2 celle);
(ii) la probabilit\`a di un \emph{ciclo} (fallimento $\Rightarrow$ rehash) si maggiora con
$\frac{1}{c-1}$: con $c=3$ \`e $\le\frac12$, quindi $\OO{1}$ rehash attesi, ammortizzato $\OO{1}$.
\emph{Prezzo}: $r\ge2cn=6n$ celle, cio\`e \kw{load factor $\le1/6\approx17\%$}: il cuckoo classico
garantisce i bound solo a tabelle molto scariche (varianti con pi\`u hash o celle multi-slot
arrivano oltre il 90\%).

\subsection*{\qid{5.4}--\qid{5.5}\hot\ Bloom Filter: definizione, bit a 0, falso positivo, ottimo}

\begin{domanda}
\qid{5.4} Calcola la probabilit\`a che una posizione del bit-array del Bloom Filter (taglia $m$,
$n$ chiavi inserite) sia 0, e deriva/dimostra la probabilit\`a d'errore complessiva. Variante:
dimostra che tale probabilit\`a \`e $\tfrac12$ quando il numero di funzioni hash \`e quello
ottimo. \emph{(11/11/20, 27/05/24, 11/11/24, 15/07/25)}\\
\qid{5.5} Definisci un Bloom Filter su un dizionario di $n$ chiavi e le operazioni supportate;
enuncia e dimostra il risultato sul falso positivo. \emph{(24/07/23)}
\end{domanda}

\textbf{Definizione e operazioni (\qid{5.5}).} Un Bloom Filter per un dizionario $D$ di $n$ chiavi
\`e un array binario $B[0,m-1]$, inizialmente tutto 0, pi\`u $k$ funzioni hash indipendenti
$h_1,\dots,h_k:U\to[0,m)$.
\emph{Insert}$(x)$: pone $B[h_i(x)]=1$ per $i=1..k$; costo $\OO{k}$.
\emph{Query}$(y)$: risponde ``presente'' sse $B[h_i(y)]=1$ \emph{per tutti} gli $i$; costo
$\OO{k}$. Non supporta la cancellazione (creerebbe falsi negativi; la variante counting/spectral
usa contatori). Propriet\`a fondamentale: \kw{niente falsi negativi, possibili falsi positivi}.

\textbf{Probabilit\`a che un bit sia 0.} Gli inserimenti accendono $kn$ posizioni casuali
indipendenti; una posizione fissata resta 0 se tutte le $kn$ estrazioni la evitano:
\[
\boxed{\ p_0=\Bigl(1-\frac{1}{m}\Bigr)^{kn}\;\approx\;e^{-kn/m}\ }
\]

\textbf{Probabilit\`a di falso positivo.} $y\notin D$ \`e dichiarata presente sse le sue $k$
posizioni sono tutte a 1 (assumendo indipendenza, approssimazione standard):
\[
\boxed{\ \varepsilon\;=\;(1-p_0)^{k}\;\approx\;\bigl(1-e^{-kn/m}\bigr)^{k}\ }
\]

\textbf{Numero ottimo di hash e $p_0=\tfrac12$ (\qid{5.4}).} Minimizziamo
$\varepsilon(k)=(1-e^{-kn/m})^k$. Posto \kw{$t=e^{-kn/m}$}, cio\`e $k=-\frac{m}{n}\ln t$:
\[
\ln\varepsilon=k\ln(1-t)=-\frac{m}{n}\,\ln t\,\ln(1-t).
\]
La funzione $g(t)=\ln t\,\ln(1-t)$ \`e \kw{simmetrica rispetto a $t\mapsto1-t$} e unimodale: il
punto ottimo \`e $t=\tfrac12$ (formalmente $g'(t)=\frac{\ln(1-t)}{t}-\frac{\ln t}{1-t}=0$ in
$t=\frac12$). Quindi
\[
e^{-k^{*}n/m}=\tfrac12
\;\Longrightarrow\;
\boxed{\;k^{*}=\frac{m}{n}\ln2\,,\qquad p_0=\tfrac12\;}
\]
\kw{con $k$ ottimo met\`a dei bit \`e a 0 e met\`a a 1} (massima ``entropia'' del filtro).
Errore all'ottimo:
\[
\boxed{\ \varepsilon^{*}=\Bigl(\tfrac12\Bigr)^{k^{*}}=2^{-\frac{m}{n}\ln2}
\approx(0.6185)^{\,m/n}\ }
\]
\kw{esponenzialmente piccolo nel numero di bit per chiave $m/n$}
(es.\ $m/n=10\Rightarrow\varepsilon\approx0.8\%$ con $k^*\approx7$).

\subsection*{\qid{5.6} Intersezione di due insiemi su due server con un Bloom Filter}

\begin{domanda}
Dati due insiemi $A$, $B$ memorizzati su due server, mostra come calcolarne l'intersezione usando
un Bloom Filter e un solo round di comunicazione, ammettendo errori. \emph{(05/06/23)}
\end{domanda}

\emph{Protocollo.} Il server 1 costruisce un Bloom Filter $BF_A$ di $A$ (taglia $m=\Theta(|A|)$
bit, $k$ ottimo) e \kw{lo invia al server 2} (un solo messaggio). Il server 2 esegue
$\textit{Query}_{BF_A}(b)$ per ogni $b\in B$ e restituisce
$X=\{b\in B:\ BF_A(b)=\text{presente}\}$.

\emph{Correttezza/errore.} \kw{Niente falsi negativi $\Rightarrow$ $A\cap B\subseteq X$}: nessun
elemento dell'intersezione \`e perso. Ogni $b\in B\setminus A$ finisce in $X$ con probabilit\`a
$\varepsilon\approx(0.6185)^{m/|A|}$: attesi $\varepsilon\,|B\setminus A|$ falsi positivi,
controllabili aumentando $m$. \emph{Comunicazione}: $m=\OO{|A|}$ \kw{bit} invece di
$\Theta(|A|\log u)$ bit per spedire $A$ esplicitamente. (Se serve l'intersezione esatta, un
secondo round di verifica elimina gli errori --- ma la domanda chiede la versione a un round.)

\subsection*{\qid{5.7} Skip List: altezza media $\OO{\log n}$}

\begin{domanda}
Dimostra che l'altezza (profondit\`a) media di una Skip List con $n$ nodi \`e $\OO{\log n}$.
\emph{(11/11/24, 04/02/25)}
\end{domanda}

\textbf{Setup.} Ogni elemento viene promosso di livello con probabilit\`a $\tfrac12$
indipendentemente: $\Pr[\text{liv}(x)\ge\ell]=2^{-(\ell-1)}$ (geometrica). L'\emph{altezza} \`e
$H=\max_x \text{liv}(x)$.

\textbf{Prova.} Per union bound, per ogni $\ell\ge1$:
\[
\Pr[H\ge\ell]\;\le\;\sum_{x}\Pr[\text{liv}(x)\ge\ell]\;=\;n\,2^{-(\ell-1)}
\qquad(\text{e ovviamente }\le1).
\]
Con l'identit\`a \kw{$\mathbb{E}[H]=\sum_{\ell\ge1}\Pr[H\ge\ell]$}, spezzando la somma in
$\ell\le1+\log_2 n$ e $\ell>1+\log_2 n$:
\[
\mathbb{E}[H]\;\le\;\underbrace{\sum_{\ell=1}^{1+\log_2 n}1}_{\le\,1+\log_2 n}
\;+\;\sum_{j\ge1} n\,2^{-(\log_2 n+j)}
\;=\;1+\log_2 n+\sum_{j\ge1}2^{-j}
\;=\;\boxed{\ \log_2 n+2\ }\;=\;\OO{\log n}. \;\blacksquare
\]
Concentrazione: $\Pr[H\ge c\log_2 n]\le 2n^{1-c}$, cio\`e \kw{altezza $\OO{\log n}$ con alta
probabilit\`a}.

\textbf{Complemento utile (ricerca $\OO{\log n}$ attesa).} Analizzando il cammino di ricerca
\emph{a ritroso}: da ogni nodo si sale di livello con probabilit\`a $\tfrac12$, altrimenti ci si
sposta orizzontalmente; per salire gli $\OO{\log n}$ livelli servono in attesa 2 passi per
livello: $\OO{\log n}$ passi totali.

\subsection*{\qid{5.8}\hot\ Treap: definizione, operazioni, split}

\begin{domanda}
Definisci la struttura Treap su $n$ coppie $\langle$chiave, priorit\`a$\rangle$ con priorit\`a
minima nella radice; descrivi le operazioni supportate e come si implementa lo split su una chiave
$K$ non presente. \emph{(10/02/23, 23/06/25, orale 25/07/22)}
\end{domanda}

\textbf{Definizione.} Un Treap su $n$ coppie $\langle k,\pi\rangle$ \`e un albero binario che \`e
\kw{contemporaneamente}: (i) un \kw{BST rispetto alle chiavi} $k$ (visita in order $=$ ordine
crescente delle chiavi); (ii) uno \kw{heap rispetto alle priorit\`a} $\pi$ --- qui un MIN-heap:
la radice ha la priorit\`a minima. \kw{Se le priorit\`a sono distinte, il Treap \`e unico}: la
radice \`e forzata (priorit\`a minima), le chiavi minori vanno a sinistra, le maggiori a destra,
ricorsivamente. Operazioni: search ($\OO{h}$ come BST), insert e delete tramite rotazioni che
ripristinano lo heap ($\OO{h}$), split e merge ($\OO{h}$); $h=\OO{\log n}$ attesa con priorit\`a
casuali (\qid{5.9}).

\textbf{Split su chiave $K$ non presente.} Obiettivo: dividere $T$ in $T_{<K}$ e $T_{>K}$.
\begin{enumerate}[leftmargin=1.8em,itemsep=1pt]
\item \kw{Si inserisce il nodo fittizio $\langle K,-\infty\rangle$} con l'insert standard: discesa
BST fino a una foglia, poi --- avendo priorit\`a minima di tutti --- il nodo \kw{risale con
rotazioni fino a diventare la radice}.
\item Per la propriet\`a BST, il sottoalbero sinistro della nuova radice contiene esattamente le
chiavi $<K$ e il destro le chiavi $>K$ (le rotazioni preservano BST e heap nei sottoalberi).
\item Si rimuove il nodo fittizio: i suoi due figli sono $T_{<K}$ e $T_{>K}$.
\end{enumerate}
Costo: un insert $=\OO{\log n}$ atteso. (Il \emph{merge} \`e l'operazione inversa: radice fittizia
sopra i due treap, che poi si fa scendere ed eliminare.)

\subsection*{\qid{5.9} Treap: profondit\`a media $\OO{\log n}$}

\begin{domanda}
Enuncia e dimostra che la profondit\`a media di un Treap con priorit\`a scelte uniformemente a
caso \`e logaritmica nel numero di chiavi. \emph{(07/02/20, orale 25/07/22)}
\end{domanda}

\textbf{Teorema.} Se le priorit\`a sono uniformi e indipendenti, per ogni nodo
$\mathbb{E}[\text{depth}(u_j)]\le 2\ln n=\OO{\log n}$.

\textbf{Lemma (struttura degli antenati).} Siano $u_1,\dots,u_n$ i nodi in ordine di chiave.
Allora \kw{$u_i$ \`e antenato di $u_j$ $\iff$ $u_i$ ha la priorit\`a minima nell'intervallo di
chiavi} $S_{ij}=\{u_{\min(i,j)},\dots,u_{\max(i,j)}\}$.

\emph{Prova del lemma.} Sia $z$ il nodo di priorit\`a minima in $S_{ij}$. Nel (unico) Treap, il
pi\`u piccolo sottoalbero che contiene tutto $S_{ij}$ ha come radice un nodo di $S_{ij}$ (le
chiavi di $S_{ij}$ sono contigue, e la radice che le separa ha chiave interna all'intervallo);
per la propriet\`a di heap tale radice \`e $z$, che quindi \`e antenato di tutti i nodi di
$S_{ij}$. Se $z=u_i$: $u_i$ antenato di $u_j$. Se $z\ne u_i$: o $z=u_j$ (e allora \`e $u_j$
l'antenato), oppure $u_i$ e $u_j$ stanno in sottoalberi \emph{diversi} di $z$ (chiavi da parti
opposte della chiave di $z$) e $u_i$ non \`e antenato di $u_j$. In tutti i casi: antenato $\iff$
minimo. $\square$

\emph{Prova del teorema.} Con priorit\`a uniformi, ogni nodo di $S_{ij}$ ($|i-j|+1$ elementi) \`e
il minimo con la stessa probabilit\`a:
\[
\Pr[u_i\text{ antenato di }u_j]=\frac{1}{|i-j|+1}.
\]
La profondit\`a di $u_j$ \`e il numero dei suoi antenati; \kw{per linearit\`a dell'attesa}:
\[
\mathbb{E}[\text{depth}(u_j)]
=\sum_{i\ne j}\frac{1}{|i-j|+1}
=\sum_{t=2}^{j}\frac{1}{t}+\sum_{t=2}^{\,n-j+1}\frac{1}{t}
\;\le\;2(H_n-1)\;\le\;\boxed{\ 2\ln n\ }=\OO{\log n}. \;\blacksquare
\]

\textbf{Commento.} \kw{Il Treap con priorit\`a casuali \`e un random BST}, qualunque sia l'ordine
di inserimento delle chiavi: \`e la sua ragion d'essere.

% ==================================================================
\section{Prefix Search --- risoluzioni}

\subsection*{\qid{6.1} Ricerca lessicografica in un Patricia Trie}

\begin{domanda}
Dato un Patricia Trie su un dizionario $S$ di $n$ stringhe, descrivi come cercare la posizione
lessicografica di $P[1,p]$ in $S$; enuncia la complessit\`a in tempo e in I/O. \emph{(24/07/23)}
\end{domanda}

\textbf{Struttura (richiamo).} Il Patricia Trie \`e il compacted trie in cui ogni arco memorizza
soltanto il \emph{primo carattere} dell'etichetta, e ogni nodo la lunghezza (skip) del prefisso
raggiunto; le stringhe complete risiedono solo nelle foglie (tipicamente su disco).

\textbf{Ricerca (blind search), tre passi:}
\begin{enumerate}[leftmargin=1.8em,itemsep=1pt]
\item \emph{Discesa cieca.} Dalla radice, a ogni nodo con skip $\ell$ si usa $P[\ell+1]$ per
scegliere l'arco (confrontando solo il carattere di diramazione, \kw{senza verificare i caratteri
saltati}). Si termina a una foglia candidata $f$ (in caso di mismatch di diramazione si prosegue
comunque su un rappresentante del sottoalbero).
\item \emph{Verifica.} Si carica \kw{l'unica stringa} $s_f$ della foglia e si calcola
$\ell^{*}=\mathrm{lcp}(P,s_f)$. Propriet\`a chiave che rende corretta la blind search:
\kw{$s_f$ massimizza l'lcp con $P$ fra tutte le stringhe di $S$} (le stringhe con lcp maggiore
starebbero sullo stesso cammino di discesa).
\item \emph{Riposizionamento.} Si risale al punto di profondit\`a $\ell^{*}$ del cammino;
confrontando $P[\ell^{*}+1]$ con i caratteri di diramazione si determina la posizione di $P$ fra i
sottoalberi, e contando le foglie a sinistra la sua posizione lessicografica in $S$.
\end{enumerate}

\textbf{Complessit\`a.}
\[
\boxed{\ \OO{p}\ \text{tempo},\qquad \OO{1}\ \text{I/O}\ }
\]
(pi\`u precisamente $\OO{p/B}$ I/O per leggere la stringa candidata). \kw{Una sola stringa
toccata per qualunque ricerca}: \`e il punto di forza del Patricia.

\subsection*{\qid{6.2}--\qid{6.3} Compacted Trie vs Patricia Trie, e l'indice two-level}

\begin{domanda}
\qid{6.2} Descrivi la differenza fra Compacted Trie e Patricia Trie in termini di spazio e tempo
di query, in funzione di $n$, della lunghezza totale $N$ e della lunghezza del pattern $P$.
\emph{(06/03/25)}\\
\qid{6.3} Qual \`e la differenza principale fra Compacted Trie e Patricia Trie, e come \`e stata
sfruttata per costruire un indice two-level efficiente per un insieme di stringhe?
\emph{(03/06/25)}
\end{domanda}

\textbf{Compacted trie.} Archi etichettati con \emph{sottostringhe intere} (coppie
(puntatore, lunghezza) nelle stringhe di $S$). Nodi $\OO{n}$, ma per navigare servono le stringhe:
spazio complessivo $\TT{N}$ caratteri $+$ $\OO{n}$ puntatori. Query: confronta \emph{tutti} i
caratteri di $P$ durante la discesa: tempo $\OO{p}$, ma se le stringhe sono su disco ogni arco
pu\`o costare accessi (pi\`u di $\OO{1}$ I/O).

\textbf{Patricia trie.} Un carattere $+$ skip per arco: \kw{spazio $\OO{n}$ parole, indipendente
da $N$}; query $\OO{p}$ con \kw{un solo accesso a una stringa} per la verifica finale: $\OO{1}$
I/O (\qid{6.1}). Prezzo: i caratteri saltati non sono controllati in discesa (serve la verifica).

\textbf{Differenza principale (\qid{6.3}).} Il compacted trie richiede le etichette complete
(spazio $\propto N$); \kw{il Patricia riduce ogni etichetta a 1 carattere $+$ 1 intero: spazio
indipendente dalla lunghezza delle stringhe}, al costo di una verifica posticipata.

\textbf{Sfruttamento nel two-level index.} Dizionario ordinato su disco a blocchi; in memoria un
\kw{Patricia trie su una stringa campione per blocco} (es.\ la prima):
\begin{itemize}[leftmargin=1.6em,itemsep=1pt]
\item lo spazio in memoria \`e $\OO{n/B_{\mathrm{blk}}}$ parole e \kw{non dipende dalla lunghezza
delle stringhe campione} --- la propriet\`a sfruttata \`e esattamente questa;
\item una query fa blind search in $\OO{p}$ tempo e $\OO{1}$ I/O, individuando \kw{l'unico blocco}
che pu\`o contenere $P$;
\item si legge quel blocco ($\OO{1}$ I/O $+$ $\OO{occ/B}$ per output lunghi) e lo si scandisce;
i blocchi sono compressi con \emph{front coding}, decodificabile dal capo-blocco memorizzato per
intero.
\end{itemize}
Risultato: prefix search in $\OO{p}$ tempo e $\OO{1+occ/B}$ I/O con memoria interna minima.

\subsection*{\qid{6.4} Storage di un dizionario di stringhe: tre soluzioni a confronto}

\begin{domanda}
Dato un dizionario $D$ di $n$ stringhe di lunghezza variabile e lunghezza totale $N$, discuti
almeno 3 soluzioni di memorizzazione, commentando spazio e costo tempo/I/O di
$\texttt{Access}(i)$. \emph{(27/05/24)}
\end{domanda}

\textbf{(1) Concatenazione $+$ array di offset.} Stringhe concatenate in un array di $N$
caratteri; array $\mathit{off}[1,n]$ ($n\lceil\log_2 N\rceil$ bit).
\emph{Access}$(i)$: lettura diretta di $[\mathit{off}[i],\mathit{off}[i+1])$:
\kw{tempo $\OO{|s_i|}$, I/O $\OO{1+|s_i|/B}$} --- ottimo. Spazio: $N$ caratteri $+$ $n\log N$ bit
(gli offset possono dominare con stringhe corte). Nessuna compressione.

\textbf{(2) Concatenazione con lunghezze in linea (senza puntatori).} Ogni stringa preceduta dalla
propria lunghezza (es.\ in $\gamma$-code): spazio $\approx N+\OO{n\log\frac{N}{n}}$ bit.
\emph{Access}$(i)$ richiederebbe la scansione dall'inizio: si corregge \kw{campionando un offset
esplicito ogni $k$ stringhe} (spazio extra $\frac{n}{k}\log N$ bit): Access $=$ salto al campione
$+$ scansione di $\le k$ stringhe: tempo $\OO{k\,\ell_{\max}}$, I/O $\OO{1+k\ell_{\max}/B}$.
\kw{Trade-off spazio/tempo regolato da $k$}.

\textbf{(3) Front coding a blocchi (FC$_B$).} Dizionario ordinato: stringhe consecutive
condividono prefissi lunghi. Blocchi di $k$ stringhe: la prima \kw{per intero} (copia di
ripartenza), le successive come coppie $\langle\ell_j,\hat{s}_j\rangle$ (lcp con la precedente,
suffisso residuo). Spazio: $FC_B(D)\ll N$ su dizionari reali ($\approx-70\%$ sugli URL) $+$
$\frac{n}{k}\log N$ bit di puntatori. \emph{Access}$(i)$: salto al blocco $\lceil i/k\rceil$
($\OO{1}$ I/O) e decodifica in scansione: tempo $\OO{k\,\ell_{\max}}$, I/O tipicamente $\OO{1}$.
\kw{\`E la soluzione usata nei blocchi dell'indice two-level} di \qid{6.3}.

\emph{Confronto}: (1) accesso ottimo, spazio massimo; (2) spazio minore, accesso pi\`u lento
(regolabile); (3) \kw{compressione vera con accesso quasi-ottimo: il compromesso standard}.
(Quarta opzione citabile: compacted trie/Patricia $+$ stringhe su disco, se servono anche ricerche
per prefisso.)

\subsection*{\qid{6.5} Ordinare stringhe con un compacted trie a fan-out hash}

\begin{domanda}
Descrivi come usare un compacted trie (con fan-out dei nodi in hash table) per ordinare $n$
stringhe di lunghezza totale $L$; commenta tempo e spazio. \emph{(10/07/20)}
\end{domanda}

\textbf{Costruzione.} Si inseriscono le $n$ stringhe una alla volta nel compacted trie, dove il
fan-out di ogni nodo \`e una \kw{hash table} (carattere di diramazione $\to$ figlio).
L'inserimento di $s$ scende carattere per carattere: ogni passo costa $\OO{1}$ atteso (lookup
hash); in caso di mismatch a met\`a arco si spezza l'arco con $\OO{1}$ nodi nuovi. Costruzione:
\kw{$\OO{L}$ atteso}; spazio $\OO{n}$ nodi $+$ etichette come puntatori nelle stringhe originali.

\textbf{Ordinamento.} L'ordine lessicografico \`e la \kw{DFS del trie con i figli visitati in
ordine alfabetico}. Ma \kw{le hash table non mantengono l'ordine}: prima della DFS si ordinano le
liste dei figli di ciascun nodo: $\sum_v \OO{d_v\log d_v}\le\OO{n\log\sigma}$ (poich\'e
$\sum_v d_v\le 2n$ e $d_v\le\sigma$). DFS finale: $\OO{n}$.

\textbf{Bilancio.}
\[
\boxed{\ \OO{L+n\log\sigma}\ \text{tempo atteso},\qquad \OO{n}\ \text{spazio extra}\ }
\]
Confronto col lower bound \qid{2.2} $\Om{d+n\log n}$: il trie evita di riconfrontare i prefissi
comuni e sposta il termine di ordinamento a $n\log\sigma$; il prezzo \`e il bound solo
\emph{atteso} (hashing) e la localit\`a peggiore rispetto al Multikey QuickSort.

% ==================================================================
\section{Substring Search --- risoluzioni}

\subsection*{\qid{7.1} Definizioni: Suffix Array e array LCP}

\begin{domanda}
Definisci formalmente il suffix array $SA$ di un testo $T[1,n]$ e il corrispondente array LCP.
\emph{(16/01/23, 02/11/23)}
\end{domanda}

\textbf{Suffix Array.} Dato $T[1,n]$ (terminato da $\$$, minimo lessicografico), $\mathit{SA}[1,n]$
\`e \kw{la permutazione di $\{1,\dots,n\}$ che ordina lessicograficamente i suffissi}:
\[
T[\mathit{SA}[1],n]\;<\;T[\mathit{SA}[2],n]\;<\;\dots\;<\;T[\mathit{SA}[n],n],
\]
cio\`e $\mathit{SA}[i]$ \`e la posizione iniziale dell'$i$-esimo suffisso in ordine
lessicografico. Spazio: $n$ interi $=\Theta(n\log n)$ bit (oltre a $T$).

\textbf{Array LCP.} $\mathit{LCP}[1,n-1]$ dove
\[
\mathit{LCP}[i]\;=\;\bigl|\mathrm{lcp}\bigl(T[\mathit{SA}[i],n],\;T[\mathit{SA}[i+1],n]\bigr)\bigr|
\]
\`e la lunghezza del pi\`u lungo prefisso comune fra i suffissi \kw{lessicograficamente
consecutivi}. Propriet\`a utile: l'lcp fra i suffissi in posizioni $i<j$ \`e
$\min_{i\le t<j}\mathit{LCP}[t]$ (range minimum).

\subsection*{\qid{7.2}, \qid{7.4}, \qid{7.5}\hot\ Ricerca e COUNT con il Suffix Array}

\begin{domanda}
\qid{7.2} Definisci $SA$ e LCP; descrivi come cercare $P[1,p]$ in $SA$; enuncia e dimostra la
complessit\`a della ricerca. \emph{(07/09/23)}\\
\qid{7.4} Mostra come CONTARE in $T[1,n]$ tutte le occorrenze di $P[1,p]$ con un Suffix Array in
memoria; enuncia e dimostra la complessit\`a; qual \`e il costo I/O del RETRIEVAL delle posizioni
se $SA$ \`e su disco? \emph{(10/02/21, 15/01/24)}\\
\qid{7.5} Descrivi come cercare $P[1,p]$ in un $SA$ costruito su $T[1,n]$ e commenta la
complessit\`a in funzione di $p$ e $n$. \emph{(15/07/25)}
\end{domanda}

\textbf{Osservazione chiave.} Le occorrenze di $P$ sono i suffissi che hanno $P$ come prefisso;
nell'ordine lessicografico questi suffissi sono \kw{consecutivi}: formano un intervallo
$\mathit{SA}[i^{*},j^{*}]$.

\textbf{Algoritmo.} \kw{Due ricerche binarie} su $\mathit{SA}$: una per il bordo sinistro $i^{*}$,
una per il bordo destro $j^{*}$. Ogni passo confronta $P$ col suffisso corrente: $\le p$ confronti
di caratteri.

\textbf{Complessit\`a (con prova).}
\[
\boxed{\ \OO{p\log n}\ \text{(ricerca ingenua)},\qquad \OO{p+\log n}\ \text{con Llcp/Rlcp}\ }
\]
\emph{Prova del primo bound}: la ricerca binaria fa $\lceil\log_2 n\rceil$ iterazioni; ogni
iterazione \`e un confronto lessicografico che termina al primo mismatch, dopo al pi\`u $p$
caratteri. \emph{Raffinamento}: mantenendo per l'intervallo corrente $[L,R]$ i valori
$\ell=\mathrm{lcp}(P,\text{suff}_{SA[L]})$, $r=\mathrm{lcp}(P,\text{suff}_{SA[R]})$ e usando gli
array precalcolati $\mathit{Llcp}/\mathit{Rlcp}$, a ogni passo \kw{o si dimezza l'intervallo senza
leggere caratteri, o ogni carattere letto di $P$ fa crescere $\max(\ell,r)$} ($\le p$ volte in
tutto): totale $\OO{p+\log n}$. $\blacksquare$

\textbf{COUNT (\qid{7.4}).} Trovato l'intervallo:
$\boxed{\mathit{occ}=j^{*}-i^{*}+1}$, senza lavoro aggiuntivo. Costo: quello della ricerca.

\textbf{RETRIEVAL con $SA$ su disco (\qid{7.4}).} Le posizioni sono
$\mathit{SA}[i^{*}..j^{*}]$: \kw{celle contigue su disco}, lette con scansione sequenziale in
\[
\boxed{\ \OO{\frac{\mathit{occ}}{B}}\ \text{I/O}\ }
\]
(pi\`u $\OO{\frac{p}{B}\log n}$ I/O per le binarie se anche $T$ \`e su disco). \kw{La contiguit\`a
rende il retrieval I/O-ottimale, output-sensitive.}

\subsection*{\qid{7.3} Costruzione del Suffix Array via qsort}

\begin{domanda}
Definisci $SA$; scrivi lo pseudocodice per costruirlo via qsort; enuncia e commenta la
complessit\`a worst-case (tempo e I/O) di questa costruzione, e il caso in cui il massimo LCP fra
i suffissi \`e limitato da una costante $c$. \emph{(07/09/21, 24/07/23)}
\end{domanda}

\textbf{Pseudocodice.}
\begin{verbatim}
BuildSA-qsort(T[1,n]):
  SA = [1, 2, ..., n]                      // indici dei suffissi
  qsort(SA, cmp)                           // comparatore:
     cmp(i, j) = confronto lessicografico  //   strncmp(T+i, T+j)
                 di T[i,n] con T[j,n]      //   carattere per carattere
  return SA
\end{verbatim}

\textbf{Tempo (worst case).} qsort fa $\OO{n\log n}$ confronti; ogni confronto costa fino a
$\Theta(n)$ caratteri (suffissi con lcp lungo, caso estremo $T=a^n$):
\[
\boxed{\ \OO{n^2\log n}\ \text{tempo worst case}\ }
\]

\textbf{I/O e perch\'e \`e inefficiente.} Due ragioni strutturali:
(a) ogni confronto scandisce fino a $\Theta(n)$ caratteri \kw{a partire da due posizioni
arbitrarie del testo}: $\OO{n/B}$ I/O a confronto, $\OO{\frac{n}{B}n\log n}$ totali;
(b) qsort permuta \emph{puntatori}: \kw{gli accessi ai suffissi sono sparsi e distruggono la
localit\`a} (niente prefetching). \`E l'esempio da manuale di algoritmo RAM-accettabile ma
I/O-pessimo; le costruzioni serie (skew/DC3, SA-IS) ottengono $\OO{n}$ o $\OO{sort(n)}$.

\textbf{Caso LCP massimo $\le c$.} Ogni confronto termina entro $c+1$ caratteri, cio\`e $\OO{1}$:
\[
\boxed{\ \OO{n\log n}\ \text{tempo}\ }
\]
\kw{Il degrado dipende dalla ripetitivit\`a del testo, misurata dagli LCP.}

\subsection*{\qid{7.6} Suffix tree: la pi\`u lunga sottostringa che occorre $\ge L$ volte}

\begin{domanda}
Dato il suffix tree di $T[1,n]$, come si calcola la sottostringa pi\`u lunga che occorre almeno
$L$ volte? Complessit\`a in funzione di $n$. \emph{(12/06/20)}
\end{domanda}

\textbf{Fatti usati.} (i) ogni sottostringa di $T$ \`e un cammino dalla radice (anche a met\`a
arco); (ii) \kw{il numero di occorrenze di una sottostringa $=$ numero di foglie nel sottoalbero}
del punto che la rappresenta; (iii) i nodi interni sono $\le n-1$.

\textbf{Algoritmo.}
\begin{enumerate}[leftmargin=1.8em,itemsep=1pt]
\item Costruisci il suffix tree (dato, oppure $\OO{n}$ con McCreight/Ukkonen).
\item Con una DFS annota ogni nodo interno $v$ con $\mathit{leaves}(v)$ (somma sui figli, foglie
$=1$): $\OO{n}$.
\item Fra i nodi interni con $\mathit{leaves}(v)\ge L$, prendi quello di \kw{massima string
depth}: la stringa radice$\to v^{*}$ \`e la risposta.
\end{enumerate}

\textbf{Correttezza.} A met\`a arco il sottoalbero (quindi il numero di occorrenze) \`e lo stesso
del nodo sottostante: la risposta ottima pu\`o essere estesa fino al prossimo nodo interno senza
perdere occorrenze, quindi \kw{termina esattamente in un nodo interno}, che l'algoritmo trova.
\textbf{Complessit\`a}: $\boxed{\OO{n}}$, indipendente da $L$.

\subsection*{\qid{7.7} Array LCP in tempo lineare (algoritmo di Kasai)}

\begin{domanda}
Scrivi lo pseudocodice dell'algoritmo lineare (Kasai et al.) che calcola l'array LCP dato $T$ e
$SA$, e dimostra che \`e $\OO{n}$. \emph{(orale 25/07/22, domanda ``da lode''; simulazioni in
11/11/24, 06/03/25, 15/07/25)}
\end{domanda}

Serve l'inverso $\mathit{ISA}$ ($\mathit{ISA}[\mathit{SA}[i]]=i$), calcolabile in $\OO{n}$.

\begin{verbatim}
Kasai(T, SA):
  for i = 1..n: ISA[SA[i]] = i        // inverso di SA
  h = 0
  for i = 1..n:                       // i = posizione NEL TESTO
      if ISA[i] > 1:
           j = SA[ ISA[i] - 1 ]       // suffisso che precede T[i,n] in ordine less.
           while i+h <= n and j+h <= n and T[i+h] == T[j+h]:
                h = h + 1             // estende il prefisso comune
           LCP[ ISA[i] - 1 ] = h
           if h > 0: h = h - 1        // proprieta' chiave (vedi prova)
      else:
           h = 0                      // T[i,n] e' il minimo lessicografico
  return LCP
\end{verbatim}

\textbf{Idea e correttezza.} Si processano i suffissi \kw{in ordine di posizione nel testo}, non
lessicografico, riusando il lavoro: se $T[i,n]$ ha lcp $=h$ col suo predecessore lessicografico
$T[j,n]$, togliendo il primo carattere si ottengono $T[i+1,n]$ e $T[j+1,n]$ con prefisso comune
$h-1$, e $T[j+1,n]$ \`e ancora minore di $T[i+1,n]$. Il predecessore \emph{immediato} di
$T[i+1,n]$ \`e, fra i minori, quello con lcp massimo: quindi \kw{il suo lcp \`e $\ge h-1$} e il
\texttt{while} del passo $i+1$ pu\`o ripartire da $h-1$ senza ricontrollare nulla.

\textbf{Complessit\`a $\OO{n}$ (prova).} Il costo dominante \`e il numero totale di incrementi di
$h$. La variabile $h$: parte da 0, \kw{non supera mai $n$, e viene decrementata al pi\`u $n$
volte} (una per iterazione del for, mai sotto 0). Quindi gli incrementi sono al pi\`u
(decrementi) $+$ (massimo) $\le 2n$:
\[
\boxed{\ \le 2n\ \text{confronti di caratteri in totale}\ \Rightarrow\ \OO{n}\ }
\]
$\blacksquare$ (\`E anche la risposta alla richiesta ``mostra il numero di caratteri
effettivamente confrontati'' delle simulazioni d'esame.)

\section{Integer Coding --- risoluzioni}

\paragraph{Convenzione sui bit alti (importante per gli esercizi).} Nei compiti del docente la
sequenza $H$ \`e costruita per \emph{bucket}: per ogni valore $v=0,1,\dots,u/2^{\ell}-1$ della
parte alta si scrivono tanti \textbf{1} quanti sono gli elementi con parte alta $=v$, seguiti da
uno \textbf{0} separatore (es.\ $H=\texttt{110\,110\,10\,0\,10}\dots$: il primo bucket ha 2
elementi, il secondo 2, il terzo 1, il quarto 0, \dots). Quindi \kw{gli 1 rappresentano gli
elementi e gli 0 i separatori}. (La dispensa usa la convenzione speculare con 0 e 1 scambiati: le
formule si trasformano scambiando $\mathrm{Select}_1\leftrightarrow\mathrm{Select}_0$.)

\subsection*{\qid{8.1}\hot\ Elias--Fano: costruzione e spazio (con prova)}

\begin{domanda}
Data una sequenza monotona di $n$ interi non negativi $<u$, enuncia e dimostra lo spazio (in bit)
della codifica di Elias--Fano. \emph{(10/01/20, 27/05/24, 04/02/25)}
\end{domanda}

\textbf{Costruzione.} Sequenza monotona $s_1\le\dots\le s_n$ di interi in $[0,u)$. Si pone
\[
\boxed{\ \ell=\Bigl\lfloor\log_2\frac{u}{n}\Bigr\rfloor\ }
\]
Ogni $s_i$ \`e spezzato in: \emph{parte bassa} $=$ gli $\ell$ bit meno significativi, scritti in
chiaro nell'array $L$ ($|L|=n\,\ell$ bit); \emph{parte alta} $h_i=\lfloor s_i/2^{\ell}\rfloor\in
[0,u/2^{\ell})$, rappresentata in $H$ con la codifica unaria a bucket.

\textbf{Teorema (spazio).}
\[
\boxed{\ |L|+|H| \;=\; n\,\ell \;+\; \Bigl(n+\frac{u}{2^{\ell}}\Bigr)
\;\;<\;\; n\Bigl\lfloor\log_2\frac{u}{n}\Bigr\rfloor+3n\ \text{bit}\ }
\]
cio\`e meno di $\log_2\frac{u}{n}+3$ bit per elemento.

\textbf{Prova.} $|L|=n\ell$: esattamente $\ell$ bit per elemento. Per $H$: \kw{un 1 per ciascuno
degli $n$ elementi e uno 0 per ciascuno dei $u/2^{\ell}$ bucket}: $|H|=n+u/2^{\ell}$. Resta da
maggiorare $u/2^{\ell}$: dalla definizione di parte intera, $\ell>\log_2\frac{u}{n}-1$, quindi
\[
2^{\ell} \;>\; \frac{u}{2n}
\qquad\Longrightarrow\qquad
\kw{\frac{u}{2^{\ell}}\;<\;2n} .
\]
Dunque $|H|<3n$ e il totale \`e $<n\ell+3n$. $\blacksquare$

\textbf{Commento.} Il minimo teorico \`e $\log_2\binom{u}{n}\approx n\log_2\frac{u}{n}+1.44n$ bit:
\kw{Elias--Fano \`e a $\approx2n$ bit dall'ottimo, qualunque sia la distribuzione}, e in cambio
supporta Access e NextGEQ efficienti (\qid{8.2}) --- ci\`o che l'interpolativo, pi\`u compatto su
sequenze clusterizzate, non offre. \kw{Negli esercizi numerici usare il conteggio esatto
$n\ell+n+u/2^{\ell}$.}

\subsection*{\qid{8.2} Access($i$) e NextGEQ($x$) su Elias--Fano}

\begin{domanda}
Scrivi l'algoritmo per eseguire $\texttt{Access}(i)$ su una sequenza codificata con Elias--Fano;
cenno a $\texttt{NextGEQ}(x)$. \emph{(11/07/24; simulazioni in molti appelli)}
\end{domanda}

\textbf{Access($i$)} (restituisce $s_i$, con $i$ da 1):
\begin{enumerate}[leftmargin=1.8em,itemsep=1pt]
\item $\mathit{pos}=\mathrm{Select}_1(H,i)$: posizione dell'$i$-esimo 1 in $H$ ($\OO{1}$ con la
struttura di select, \qid{10.1});
\item \kw{parte alta: $h_i=\mathit{pos}-i$} --- prima dell'$i$-esimo 1 ci sono $i-1$ uni e
$\mathit{pos}-i$ zeri, e il numero di separatori che precedono l'elemento \`e l'indice del suo
bucket;
\item parte bassa: $L[i]$, cio\`e i bit $[(i-1)\ell+1,\,i\ell]$ di $L$;
\item output:
\[
\boxed{\ s_i=(\mathrm{Select}_1(H,i)-i)\cdot 2^{\ell}+L[i]\ }
\]
\end{enumerate}
Costo: \kw{$\OO{1}$} (una select $+$ una lettura di $\ell$ bit consecutivi).

\textbf{NextGEQ($x$)}: sia $b=\lfloor x/2^{\ell}\rfloor$ il bucket di $x$. Con
$p=\mathrm{Select}_0(H,b)$ si trova la fine del bucket $b-1$: gli elementi del bucket $b$
iniziano dall'indice $j_0=p-b+1$. Si scandiscono le parti basse del bucket $b$ confrontandole con
$x\bmod 2^{\ell}$; il primo elemento $\ge x$ \`e la risposta; se il bucket si esaurisce, \`e il
primo elemento del bucket non vuoto successivo. Costo: $\OO{1}$ select $+$ scansione del bucket
(ammortizzata $\OO{1}$ nell'uso tipico, \qid{4.4}).

\subsection*{\qid{8.3} Ricavare $n$ (e $\ell$) dagli array $L$, $H$}

\begin{domanda}
Dati gli array $L$ e $H$ di una codifica Elias--Fano, mostra come ricavare il valore di $n$
(numero di interi codificati), motivando la risposta. \emph{(orale 25/07/22; \`e il ``hint''
ricorrente degli esercizi: 11/11/20, 12/12/23, 09/12/24, 06/03/25, 23/06/25)}
\end{domanda}

\textbf{Risposta.}
\[
\boxed{\ n=\#_1(H)\ \text{ (numero di bit a 1 in $H$)}\ }
\]

\textbf{Motivazione.} Per costruzione di $H$, \kw{ogni elemento della sequenza contribuisce con
esattamente un bit a 1} (nel bucket della sua parte alta), e gli 0 sono solo separatori:
$\#_1(H)=n$ indipendentemente dalla distribuzione nei bucket. Da qui si ricava tutto:
\[
\kw{\ell=\frac{|L|}{n}}\quad(\text{$\ell$ bit per elemento in }L),
\qquad
\#_0(H)=\frac{u}{2^{\ell}}\ (\text{numero di bucket}),
\qquad
b=\ell+\log_2\#_0(H).
\]
\`E il procedimento degli hint d'esame: es.\ $|L|=22$, $\#_1(H)=11\Rightarrow n=11$, $\ell=2$;
con $\#_0(H)=16$ bucket, interi originali su $2+4=6$ bit.

\subsection*{\qid{8.4} Lunghezze dei codici (riferimento per gli esercizi di calcolo)}

\begin{domanda}
(Riferimento) Numero di bit emessi dai codici unario, $\gamma$, $\delta$, Rice e $(s,c)$-dense su
un intero $x$. \emph{(vari appelli)}
\end{domanda}

Detta $B(x)$ la rappresentazione binaria di $x\ge1$ ($|B(x)|=\lfloor\log_2 x\rfloor+1$ bit):
\begin{itemize}[leftmargin=1.6em,itemsep=2pt]
\item \textbf{Unario} $U(x)$: \kw{$x$ bit}.
\item \textbf{$\gamma$}: $|B(x)|-1$ zeri seguiti da $B(x)$:
\kw{$|\gamma(x)|=2\lfloor\log_2 x\rfloor+1$}. Es.: $\gamma(9)=\texttt{000\,1001}$, 7 bit.
\item \textbf{$\delta$}: $\gamma(|B(x)|)$ seguito da $B(x)$ \emph{intero} (convenzione della
dispensa): $|\delta(x)|\approx\log_2x+2\log_2\log_2x$.
Es.: $x=14$: $B(14)=\texttt{1110}$, $\gamma(4)=\texttt{00\,100}$,
$\delta(14)=\texttt{00\,100\,1110}$, \kw{9 bit}. (Alcuni testi omettono il bit pi\`u
significativo di $B(x)$: all'esame usare la convenzione della dispensa.)
\item \textbf{Rice $R_k$}: quoziente $q=\bigl\lfloor\frac{x-1}{2^{k}}\bigr\rfloor$ in unario
($q+1$ bit) $+$ resto su $k$ bit: \kw{$|R_k(x)|=\lfloor(x-1)/2^{k}\rfloor+k+1$}.
Es.: $R_4(83)$: $q=\lfloor82/16\rfloor=5$, $5+4+1=10$ bit. Ottimo per interi concentrati intorno
a $2^{k}$.
\item \textbf{$(s,c)$-dense}: alfabeto di $2^{b}$ simboli ($b$ bit l'uno), $s$ \emph{stopper} $+$
$c$ \emph{continuer} ($s+c=2^{b}$): codeword $=$ sequenza di continuer chiusa da uno stopper;
\kw{$s$ codeword da 1 simbolo, $cs$ da 2, $c^{2}s$ da 3}, ecc. (negli esercizi: enumerare).
\item \textbf{Elias--Fano}: $\approx\log_2\frac{u}{n}+3$ bit per intero; esatto:
$n\ell+n+u/2^{\ell}$ (\qid{8.1}).
\end{itemize}

% ==================================================================
\section{Statistical Coding --- risoluzioni}

\subsection*{\qid{9.1}\hot\ Canonical Huffman: perch\'e, preambolo, decodifica}

\begin{domanda}
Descrivi perch\'e \`e stato introdotto il Canonical Huffman; specifica quali strutture dati tiene
nel preambolo del file compresso; scrivi lo pseudocodice per decomprimere un simbolo.
\emph{(09/02/24, 21/06/24, 15/07/25)}
\end{domanda}

\textbf{Perch\'e si introduce.} L'albero di Huffman esplicito ha due difetti:
(i) \emph{spazio}: memorizzarne la struttura nel preambolo costa $\Theta(\sigma)$ puntatori;
(ii) \emph{tempo}: decodificare seguendo puntatori radice$\to$foglia fa \kw{un salto di memoria
per bit} (cache miss). Osservazione risolutiva: \kw{dell'albero contano solo le lunghezze
$L(\sigma)$} --- qualunque assegnamento prefix-free con le stesse lunghezze \`e ottimo (Kraft). Il
canonico assegna ai simboli di lunghezza $\ell$, ordinati, \kw{valori binari consecutivi}: la
decodifica lavora con due array e confronti aritmetici (al pi\`u 2 cache miss per simbolo).

\textbf{Preambolo.} Per ogni lunghezza $\ell=1,\dots,\ell_{\max}$:
\begin{itemize}[leftmargin=1.6em,itemsep=1pt]
\item $\mathit{num}[\ell]$: quanti simboli hanno codeword di lunghezza $\ell$;
\item $\mathit{fc}[\ell]$ (\emph{first code}): valore della prima codeword di lunghezza $\ell$,
con la ricorrenza \kw{dal livello pi\`u profondo verso l'alto}
\[
\boxed{\ \mathit{fc}[\ell_{\max}]=0,\qquad
\mathit{fc}[\ell]=\Bigl\lceil\frac{\mathit{fc}[\ell+1]+\mathit{num}[\ell+1]}{2}\Bigr\rceil\ }
\]
\item $\mathit{symb}[\ell]$: lista dei simboli di lunghezza $\ell$; la codeword del $t$-esimo \`e
il numero $\mathit{fc}[\ell]+t-1$ su $\ell$ bit.
\end{itemize}
(Sanity check con le tabelle d'esame: $\mathit{num}=[0,3,1,2]$ d\`a $\mathit{fc}[4]=0$,
$\mathit{fc}[3]=1$, $\mathit{fc}[2]=1$, $\mathit{fc}[1]=2$, cio\`e $FC=\{2,1,1,0\}$;
\kw{$\mathit{fc}[1]=2$ non \`e rappresentabile su 1 bit: \`e la sentinella} ``nessuna codeword di
lunghezza 1''.)

\textbf{Decodifica di un simbolo.} Propriet\`a: una sequenza di bit letta come numero $v$ su
$\ell$ bit \`e una codeword completa \kw{sse $v\ge \mathit{fc}[\ell]$}.
\begin{verbatim}
DecodeSymbol():
  v = next_bit();  len = 1
  while v < fc[len]:                 // non e' ancora una codeword completa
      v = 2*v + next_bit()           // aggiungi un bit
      len = len + 1
  return symb[len][ v - fc[len] ]    // offset dentro il livello
\end{verbatim}
(Esempio d'esame 13/06/22: $FC=\{2,1,1,0\}$, $\mathit{symb}=\{[\,],[a,e,f],[c],[d,b]\}$, input
\texttt{10001}: $v=1<2$; $v=10_2=2\ge1$ $\Rightarrow$ $symb[2][1]=e$; poi \texttt{001}: $v=0<2$,
$v=0<1$, $v=1\ge1$ $\Rightarrow$ $symb[3][0]=c$.)

\subsection*{\qid{9.2} Upper bound in bit della codifica aritmetica}

\begin{domanda}
Dimostra l'upper bound in bit della codifica aritmetica, in funzione dell'entropia e della
lunghezza del testo in input. \emph{(12/12/23, 09/12/24)}
\end{domanda}

\textbf{Teorema.} Sia $T=s_1\cdots s_n$ compresso con codifica aritmetica; se $\mathcal{P}$ \`e la
distribuzione empirica di $T$ ed $H$ l'entropia empirica (ordine 0), i bit emessi sono
\[
\boxed{\ \Bigl\lceil\log_2\frac{1}{\prod_{k}\mathcal{P}[s_k]}\Bigr\rceil+1\;\le\;nH+2\ }
\]

\textbf{Prova (tre passi).}

\emph{(1) Ampiezza dell'intervallo finale.} Da $[0,1)$, ogni simbolo $s_k$ restringe l'intervallo
di un fattore $\mathcal{P}[s_k]$: dopo $n$ simboli
\[
\kw{s\;=\;\prod_{k=1}^{n}\mathcal{P}[s_k]}.
\]

\emph{(2) Bastano $\lceil\log_2\frac1s\rceil+1$ bit.} Si emette il \kw{punto medio
$c=lo+\frac{s}{2}$ troncato} a $k=\lceil\log_2\frac{1}{s}\rceil+1$ bit frazionari. Per il bound
sul troncamento (\qid{9.3}): $c-z<2^{-k}\le\frac{s}{2}$, quindi
\[
z\;>\;c-\frac{s}{2}\;=\;lo,\qquad z\;\le\;c\;<\;lo+s :
\]
\kw{il numero emesso resta dentro l'intervallo finale}, e il decoder (che riproduce la stessa
partizione ricorsiva di $[0,1)$) recupera univocamente $T$.

\emph{(3) Dal prodotto all'entropia.} Prendendo i logaritmi:
\[
\log_2\frac{1}{s}=\sum_{k=1}^{n}\log_2\frac{1}{\mathcal{P}[s_k]}
=\sum_{\sigma\in\Sigma}n_\sigma\log_2\frac{1}{\mathcal{P}[\sigma]}
\;\overset{\mathcal{P}[\sigma]=n_\sigma/n}{=}\;n\,H .
\]
Quindi $k\le\log_2\frac1s+2=nH+2$. $\blacksquare$

\textbf{Commento.} \kw{L'overhead \`e di 2 bit sull'intero testo, non per simbolo}: \`e ci\`o che
rende l'aritmetico quasi ottimo anche con $H\ll1$, dove Huffman paga $\ge1$ bit/simbolo
(\qid{9.4}).

\subsection*{\qid{9.3} Errore di troncamento di un numero binario $<1$}

\begin{domanda}
Dimostra un upper bound alla variazione indotta dal troncamento a $k$ cifre di un numero $<1$
rappresentato in binario come $.b_1b_2b_3\ldots$ \emph{(10/01/20)}
\end{domanda}

\textbf{Teorema.} Sia $x=.b_1b_2b_3\ldots=\sum_{i\ge1}b_i2^{-i}\in[0,1)$ e $x_k=.b_1\ldots b_k$ il
troncamento alle prime $k$ cifre. Allora
\[
\boxed{\ 0\;\le\;x-x_k\;=\;\sum_{i>k}b_i\,2^{-i}\;\le\;\sum_{i>k}2^{-i}\;=\;2^{-k}\ }
\]
cio\`e \kw{il troncamento a $k$ cifre riduce il valore di al pi\`u $2^{-k}$} (serie geometrica
della coda; uguaglianza solo per coda tutta a 1). $\blacksquare$

\textbf{Uso.} \`E il lemma del passo (2) di \qid{9.2}: troncando il punto medio a $k$ cifre con
$2^{-k}\le s/2$ si resta dentro l'intervallo finale.

\subsection*{\qid{9.4} Huffman vs Arithmetic: bit emessi e quando usare l'uno o l'altro}

\begin{domanda}
Enuncia il numero di bit emessi da Huffman e da Arithmetic dato un testo $T$ di lunghezza $n$ ed
entropia empirica $H$; commenta quando conviene comprimere con l'uno o con l'altro, tenendo conto
anche dell'efficienza in tempo. \emph{(06/03/25, 03/06/25)}
\end{domanda}

\textbf{Bit emessi.}
\[
\boxed{\ \text{Huffman: } nH\;\le\;|C(T)|\;<\;n(H+1)
\qquad\quad
\text{Arithmetic: } |C(T)|\;\le\;nH+2\ }
\]
Per Huffman: la lunghezza media per simbolo dell'ottimo prefix-free soddisfa $H\le L<H+1$ (lower
bound di Shannon; upper bound perch\'e il codice di Shannon con lunghezze
$\lceil\log\frac1{p_\sigma}\rceil$ \`e prefix-free e Huffman non \`e peggio). \kw{Lo spreco di
Huffman pu\`o avvicinarsi a 1 bit/simbolo}: ogni codeword \`e lunga almeno 1, anche per simboli
con $\log\frac1p\approx0$.

\textbf{Quando conviene.}
\begin{itemize}[leftmargin=1.6em,itemsep=1pt]
\item \emph{Compressione}: \kw{l'aritmetico domina su distribuzioni sbilanciate ($H\ll1$)} ---
es.\ output di MTF/RLE0 nella pipeline bzip2: Huffman emette comunque $\ge n$ bit, l'aritmetico
$\approx nH+2\ll n$. Con $H$ grande la differenza \`e trascurabile.
\item \emph{Tempo}: \kw{Huffman canonico \`e molto pi\`u veloce} (lookup e aritmetica intera,
\qid{9.1}); l'aritmetico paga moltiplicazioni e rinormalizzazioni per simbolo.
\item \emph{Modelli adattivi/contestuali}: l'aritmetico si sposa con probabilit\`a che cambiano a
ogni simbolo (PPM); Huffman dovrebbe ricostruire il codice a ogni aggiornamento.
\end{itemize}
Sintesi: \kw{distribuzioni skewed o modelli adattivi $\Rightarrow$ aritmetico; velocit\`a e
distribuzioni non estreme $\Rightarrow$ Huffman canonico.}

\subsection*{\qid{9.5} L'aritmetico \`e invariante per permutazione del testo}

\begin{domanda}
Dimostra che due testi $T_1$, $T_2$, l'uno permutazione dell'altro, sono compressi con lo stesso
numero di bit dalla codifica aritmetica. \emph{(23/06/25)}
\end{domanda}

\textbf{Prova.} Dalla prova di \qid{9.2}, il numero di bit dipende \kw{solo dall'ampiezza
dell'intervallo finale}:
\[
\text{bit}(T)=\Bigl\lceil\log_2\frac{1}{s(T)}\Bigr\rceil+1,
\qquad
s(T)=\prod_{k=1}^{n}\mathcal{P}[t_k]
=\kw{\prod_{\sigma\in\Sigma}\mathcal{P}[\sigma]^{\,n_\sigma(T)}} .
\]
L'ultimo prodotto dipende soltanto dai \kw{conteggi} $n_\sigma(T)$, non dall'ordine dei simboli
(commutativit\`a). $T_1$ e $T_2$, essendo l'uno permutazione dell'altro, hanno gli stessi conteggi
e lo stesso modello statico $\mathcal{P}$: $s(T_1)=s(T_2)$, stesso numero di bit. $\blacksquare$

(I bit emessi possono \emph{differire come stringa} --- gli intervalli sono diversi --- ma il loro
\emph{numero} coincide. La propriet\`a vale per modelli statici e count-based; \`e falsa per
modelli contestuali di ordine $\ge1$.)

\subsection*{\qid{9.6} Perch\'e si emette il centro (troncato) e non l'estremo sinistro}

\begin{domanda}
Perch\'e la codifica aritmetica non emette l'estremo sinistro dell'intervallo finale, ma parte
dall'elemento centrale e lo tronca? \emph{(07/09/21)}
\end{domanda}

L'encoder deve trasmettere abbastanza bit da identificare \emph{un numero interno}
all'intervallo finale $[lo,\,lo+s)$. Due problemi con $lo$:
\begin{itemize}[leftmargin=1.6em,itemsep=1pt]
\item \kw{$lo$ pu\`o richiedere moltissimi bit (anche infiniti)}: la sua espansione binaria non ha
alcun legame con l'ampiezza $s$; trasmetterlo per intero vanificherebbe il bound
$\approx\log_2\frac1s$;
\item \kw{non si pu\`o troncare}: il troncamento riduce il valore (\qid{9.3}) e qualunque numero
$<lo$ \`e \emph{fuori} dall'intervallo: il decoder sbaglierebbe.
\end{itemize}
Soluzione: partire dal \kw{centro $c=lo+\frac{s}{2}$} e troncarlo a
$k=\lceil\log_2\frac1s\rceil+1$ bit: per \qid{9.3} si scende di $<2^{-k}\le\frac{s}{2}$, quindi si
resta $>lo$, dentro l'intervallo, col minimo numero di bit legato solo a $s$. \kw{Il centro \`e
l'unico punto che tollera un errore di troncamento pari a $s/2$.}

% ==================================================================
\section{Strutture dati compresse --- risoluzioni}

\subsection*{\qid{10.1}--\qid{10.2}\hot\ Rank in $\OO{1}$ tempo e $o(m)$ bit}

\begin{domanda}
\qid{10.1} Descrivi la struttura dati per supportare $\texttt{Rank}$ su un array binario $B[1,n]$
in tempo $\OO{1}$ e valutane lo spazio in bit. \emph{(16/01/23, 05/06/23, 02/11/23)}\\
\qid{10.2} Dimostra la complessit\`a in tempo e in spazio (in bit) della soluzione che implementa
$\texttt{Rank}$ in $\OO{1}$ e $o(m)$ bit su $B[1,m]$. \emph{(12/12/23, 10/01/25, 04/02/25)}
\end{domanda}

\textbf{Problema.} Dato $B[1,m]$ binario statico, supportare
$\mathrm{Rank}_1(x)=\#\{i\le x: B[i]=1\}$ in $\OO{1}$ con spazio \emph{aggiuntivo} $o(m)$ bit.
(Da $\mathrm{Rank}_1$: $\mathrm{Rank}_0(x)=x-\mathrm{Rank}_1(x)$.)

\textbf{Struttura a tre livelli.} Parametri:
\kw{$Z=(\log_2 m)^2$ (big block) e $z=\frac12\log_2 m$ (small block)}.
\begin{enumerate}[leftmargin=1.8em,itemsep=1pt]
\item \emph{Big block}: per ognuno dei $m/Z$ blocchi grandi, il rank \kw{assoluto} all'inizio del
blocco ($\log_2 m$ bit ciascuno).
\item \emph{Small block}: per ognuno dei $m/z$ blocchi piccoli, il rank \kw{relativo} all'inizio
del suo big block (valori $\le Z$: $\OO{\log\log m}$ bit ciascuno).
\item \emph{Lookup table} $R$: per ogni configurazione di $z$ bit e ogni offset $o\le z$,
$R[b,o]=$ numero di 1 nei primi $o$ bit di $b$ (condivisa da tutti i blocchi).
\end{enumerate}

\textbf{Query.}
\[
\boxed{\ \mathrm{Rank}_1(x)=\mathit{Abs}\Bigl[\tfrac{x}{Z}\Bigr]
+\mathit{Rel}\Bigl[\tfrac{x}{z}\Bigr]
+R[\,B_{\text{small}(x)},\;x\bmod z\,]\ }
\]
tre letture di array e una somma: \kw{tempo $\OO{1}$ worst case} (lo small block di
$\frac12\log m$ bit si legge in $\OO{1}$ operazioni su parola di macchina).

\textbf{Analisi dello spazio (la parte da dimostrare).}
\begin{itemize}[leftmargin=1.6em,itemsep=2pt]
\item Big block: $\dfrac{m}{Z}\log_2 m=\dfrac{m}{(\log m)^2}\log m=\kw{\dfrac{m}{\log m}=o(m)}$
--- serve $Z=(\log m)^2$ per ammortizzare i $\log m$ bit del contatore assoluto.
\item Small block: $\dfrac{m}{z}\,\OO{\log\log m}=\dfrac{2m}{\log m}\,\OO{\log\log m}
=\kw{\OO{\dfrac{m\log\log m}{\log m}}=o(m)}$ --- serve la \emph{gerarchia}: il contatore relativo
costa solo $\log Z=\OO{\log\log m}$ bit, non $\log m$.
\item Lookup table: configurazioni $2^{z}=2^{\frac12\log m}=\kw{\sqrt{m}}$; ciascuna con $z+1$
offset e conteggi da $\OO{\log z}$ bit: $\sqrt{m}\cdot\OO{\log m\log\log m}=o(m)$ --- serve
$z=\frac12\log m$: \kw{l'esponente $\frac12$ rende la tabella sub-lineare} (con blocchi di
$\log m$ bit sarebbe $\Theta(m\,\mathrm{polylog})$).
\end{itemize}
Totale:
\[
\boxed{\ \OO{\frac{m\log\log m}{\log m}}=o(m)\ \text{bit aggiuntivi},\qquad
\text{query }\OO{1}\ }
\]
$\blacksquare$

(\emph{Esercizi numerici} --- 07/09/23, 10/01/25, 23/06/25 --- con parametri dati, es.\ $Z=6$,
$z=2$: costruire i tre livelli e risolvere $\mathrm{Rank}(x)$ sommando i tre contributi.)

\textbf{Select (cenno).} $\mathrm{Select}_1(j)$ (posizione del $j$-esimo 1) si ottiene con analoga
gerarchia in $\OO{1}$ tempo e $o(m)$ bit: \`e la primitiva usata da Elias--Fano (\qid{8.2}) e
dagli alberi succinti (\qid{10.3}).

\subsection*{\qid{10.3} Alberi binari succinti e LOUDS (riferimento operativo)}

\begin{domanda}
(Riferimento) Codifica succinta di alberi binari in $B[1,2n+1]$ bit e navigazione con
Rank/Select; LOUDS per alberi generali. \emph{(base degli esercizi: 11/11/20, 15/12/21, 25/07/22,
16/01/23, 10/02/23, 24/07/23, 12/12/23, 09/02/24, 10/01/25, 04/02/25, 06/03/25, 03/06/25,
23/06/25)}
\end{domanda}

\textbf{Codifica in $2n+1$ bit.} Si completa l'albero con le foglie esterne (NULL): ogni nodo
reale ha 2 figli (reali o esterni), i nodi esterni sono $n+1$. Visita \kw{per livelli (BFS)}
scrivendo \texttt{1} per ogni nodo reale e \texttt{0} per ogni nodo esterno: $B[1,2n+1]$.
Etichette in un array $L$ in ordine BFS.

\textbf{Navigazione} (nodo $=$ posizione $x$ in $B$ con $B[x]=1$; sia $i=\mathrm{Rank}_1(x)$ il
suo numero BFS):
\[
\boxed{\ \mathrm{leftChild}(x)=2i,\qquad
\mathrm{rightChild}(x)=2i+1,\qquad
\mathrm{parent}(x)=\mathrm{Select}_1\bigl(\lfloor x/2\rfloor\bigr)\ }
\]
(ogni nodo reale genera esattamente 2 posizioni nel livello successivo). \kw{Test di foglia:
$B[2i]=0$ e $B[2i+1]=0$}. Etichetta di $x$: $L[\mathrm{Rank}_1(x)]$. Tutto $\OO{1}$ con
Rank/Select (\qid{10.1}). \`E l'attrezzo per gli esercizi ricorrenti: navigare cammini, contare
non-foglie, stampare etichette di foglie, verificare archi $(A,B)$, ecc.

\textbf{LOUDS per alberi generali.} Visita BFS scrivendo per ogni nodo il grado in unario
($\texttt{1}^{\deg}\texttt{0}$), con super-radice iniziale \texttt{10}: $2n+\OO{1}$ bit. Il
$j$-esimo nodo BFS \`e il $j$-esimo \texttt{1};
$\mathrm{child}(j,t)=\mathrm{Rank}_1(\mathrm{Select}_0(j))+t$; il padre si trova col \texttt{0}
che chiude il gruppo del $j$-esimo \texttt{1}: $\OO{1}$ con Rank/Select.

\textbf{Nota su FM-index (possibile domanda di frontiera).} La backward search mantiene
l'intervallo $[\mathit{first},\mathit{last}]$ delle righe della matrice BWT prefissate da
$P[i,p]$, aggiornandolo per $c=P[i]$ con
$\mathit{first}=C[c]+\mathrm{Rank}_c(\mathit{first}-1)+1$,
$\mathit{last}=C[c]+\mathrm{Rank}_c(\mathit{last})$: $p$ passi da $\OO{t_{\mathrm{Rank}}}$, cio\`e
\kw{$\OO{p}$ con Rank $\OO{1}$}; correttezza dall'LF-mapping della BWT.

% ==================================================================
\section{Domande di progettazione --- risoluzioni}

\subsection*{\qid{11.1} Massima somma di valori satellite (07/02/20)}

\begin{domanda}
\`E data una sequenza $S$ di $n$ coppie $\langle$stringa, occ$\rangle$, dove pi\`u coppie possono
riferirsi alla stessa stringa e occ \`e un valore satellite intero; sia $s$ il numero di stringhe
distinte. Progetta un algoritmo I/O-efficiente che calcola la stringa con massima somma dei valori
satellite, valutandone la complessit\`a I/O nei tre casi: (a) $M$ molto pi\`u piccola di $s$;
(b) $M=\OO{s\log n}$ bit ammettendo errori in probabilit\`a; (c) $M$ pu\`o contenere $\OO{s}$
puntatori a stringhe e $\OO{s}$ caratteri e contatori (hint: soluzione trie-based, ma non un trie
classico). \emph{(07/02/20)}
\end{domanda}

Sia $N$ la lunghezza totale in caratteri dei dati.

\textbf{(a) $M\ll s$.} Non si pu\`o accumulare in memoria: \kw{si ordina}. MergeSort multi-way
(\qid{1.1}) sulle coppie con chiave la stringa: $\OO{\frac{N}{B}\log_{M/B}\frac{N}{M}}$ I/O. Poi
una scansione: le coppie della stessa stringa sono adiacenti; si accumula la somma corrente e si
mantiene il massimo: $\OO{N/B}$ I/O. Totale: $\boxed{\OO{sort(N)}}$ I/O.

\textbf{(b) $M=\OO{s\log n}$ bit, errori ammessi.} \kw{Contatori hashati con fingerprint}: tabella
di $\Theta(s)$ celle, ognuna con un contatore ($\OO{\log n}$ bit) $+$ un fingerprint di
$\OO{\log s}$ bit (hash universale secondario). Per ogni coppia: $T[h(w)]\mathrel{+}=o$ se il
fingerprint coincide. Con fingerprint di $c\log_2 s$ bit, la probabilit\`a che una qualche coppia
di stringhe distinte collida su cella e fingerprint \`e $\le\binom{s}{2}/s^{c}$: trascurabile. Si
tiene copia della miglior stringa corrente. \kw{Una passata: $\OO{N/B}$ I/O}, risposta corretta
con alta probabilit\`a. (Alternativa: Count-Min sketch, con sovrastima controllata.)

\textbf{(c) $\OO{s}$ puntatori $+$ $\OO{s}$ caratteri e contatori $\Rightarrow$ \kw{Patricia
trie}.} Un trie classico costerebbe troppi caratteri in memoria; \kw{il Patricia costa $\OO{s}$
nodi con un carattere $+$ uno skip per nodo} $+$ un contatore e un puntatore su disco per foglia:
esattamente il budget. Per ogni coppia $\langle w,o\rangle$: blind search di $w$ ($\OO{|w|}$, 0
I/O); \kw{verifica sulla stringa puntata dalla foglia candidata (1 I/O)}: se uguale, contatore
$\mathrel{+}=o$; se diversa, $w$ \`e nuova: la si appende al file su disco e si inserisce una
foglia usando la posizione di mismatch appena calcolata. Costo: tempo $\OO{N}$, I/O
$\boxed{\OO{n+N/B}}$, memoria $\OO{s}$. \kw{Rispetto ad (a) evita il sort; rispetto a (b) \`e
esatta.}

\subsection*{\qid{11.2} Finestra di $k$ posizioni con $\ge q$ occorrenze di $P$ (02/02/22)}

\begin{domanda}
Dato un testo $T[1,n]$ indicizzato da un Suffix Array, progetta un algoritmo che, dati $P[1,p]$ e
due interi positivi $k$ e $q$, stabilisce se esiste un range di $k$ posizioni contigue
$T[i,i+k-1]$ in cui iniziano almeno $q$ occorrenze di $P$; stampa $i$ se esiste, altrimenti $-1$.
\emph{(02/02/22)}
\end{domanda}

\textbf{Algoritmo.}
\begin{enumerate}[leftmargin=1.8em,itemsep=1pt]
\item Ricerca di $P$ in $SA$: intervallo $SA[i^{*},j^{*}]$ delle occorrenze (\qid{7.2}):
$\OO{p\log n}$. Sia $occ=j^{*}-i^{*}+1$; \kw{se $occ<q$: stampa $-1$} e termina.
\item Estrai le posizioni $\mathit{pos}=SA[i^{*}..j^{*}]$ e \kw{ordinale} (in $SA$ sono in ordine
lessicografico, non testuale): $\OO{occ\log occ}$.
\item \kw{Sliding window}: per $t=1,\dots,occ-q+1$ controlla se
$\mathit{pos}[t+q-1]-\mathit{pos}[t]\le k-1$ ($q$ occorrenze consecutive nel testo dentro una
finestra di $k$ posizioni). Al primo $t$ valido stampa $i=\mathit{pos}[t]$; altrimenti $-1$.
$\OO{occ}$.
\end{enumerate}
\textbf{Correttezza.} Se una finestra contiene $\ge q$ inizi, fra le posizioni ordinate ce ne sono
$q$ consecutive nella finestra, e la coppia (prima, $q$-esima) dista $\le k-1$: il test la trova;
viceversa, se il test passa, la finestra che parte da $\mathit{pos}[t]$ contiene le $q$
occorrenze.
\textbf{Complessit\`a.} $\boxed{\OO{p\log n+occ\log occ}}$ tempo, $\OO{occ}$ spazio.

\subsection*{\qid{11.3} Parole $>W$ con frequenza in $[f_{\min},f_{\max}]$ (11/11/24)}

\begin{domanda}
\`E dato un insieme di coppie $\langle$frequenza, parola$\rangle$. Progetta una struttura dati e
un algoritmo che, dati un range di frequenze $[f_{\min},f_{\max}]$ e una parola $W$, restituiscono
tutte le parole lessicograficamente maggiori di $W$ con frequenza nel range; discuti
complessit\`a in tempo e spazio. \emph{(11/11/24)}
\end{domanda}

\`E una query bidimensionale (ordine lessicografico $\times$ frequenza). Soluzione nello spirito
del corso: \kw{Treap usato come priority search tree}: \emph{chiave} $=$ parola (ordine
lessicografico), \emph{priorit\`a} $=$ frequenza a MAX-heap (radice $=$ frequenza massima). Le
priorit\`a non sono casuali ma dati del problema: la struttura resta valida (BST $\times$ heap),
si perde solo la garanzia probabilistica di bilanciamento (costruzione una tantum
$\OO{n\log n}$).

\textbf{Query} $(W,f_{\min},f_{\max})$: visita ricorsiva con \kw{doppia potatura}:
\begin{verbatim}
Report(v):
  if v = NIL: return
  if freq(v) < f_min: return            // MAX-heap: tutto il sottoalbero < f_min
  if key(v) > W:
       if freq(v) <= f_max: output (key(v), freq(v))
       Report(left(v));  Report(right(v))
  else:                                  // key(v) <= W: a sinistra tutto <= W
       Report(right(v))
\end{verbatim}
Correttezza delle potature: (i) \kw{per il MAX-heap, se $freq(v)<f_{\min}$ nessun discendente \`e
in range}; (ii) \kw{per il BST, se $key(v)\le W$ tutto il sottoalbero sinistro ha chiavi $\le W$}.

\textbf{Complessit\`a.} Spazio $\OO{n}$. Query $\OO{h+occ'}$ con $h$ altezza e
$occ'=\#\{$parole $>W$ con freq $\ge f_{\min}\}$: per la query three-sided
$(>W)\times[f_{\min},\infty)$ \`e l'output-sensitive classico $\OO{h+occ}$; il vincolo $\le
f_{\max}$ non aiuta la potatura (i nodi troppo frequenti vanno comunque attraversati).
\emph{Alternativa con garanzie}: range tree con query $\OO{\log^2 n+occ}$ e spazio
$\OO{n\log n}$ --- da citare come trade-off.


\end{document}
